\documentclass[10pt,german,a4paper]{scrreprt}

\usepackage[ngerman]{babel}
\usepackage[utf8]{inputenc} 
\usepackage{lmodern}
\usepackage{listings}
\usepackage[usenames,dvipsnames]{color}
\usepackage{amsmath}
\usepackage{amssymb}
\usepackage{pdfpages}
\usepackage{geometry}
\usepackage{fancyhdr}
\usepackage{graphicx}
\usepackage{hyperref} 
\usepackage{float}		
\usepackage{relsize}		
\usepackage{euscript}	
\usepackage{titlesec}

\definecolor{battleshipgrey}{rgb}{0.52, 0.52, 0.51}

\hypersetup{
	pdfborder=0
}

\geometry{a4paper,left=1cm,right=1cm, top=2cm, bottom=2cm} 

\def \Path {/media/Daten/Studium/Skripte/TechnischeInformatik}
\makeatletter
	\@addtoreset{chapter}{part}
\makeatother

\titleformat*{\section}{\sffamily\LARGE\bfseries\color{battleshipgrey}}
\titleformat*{\subsection}{\sffamily\Large\bfseries\color{battleshipgrey}}
\titleformat*{\subsubsection}{\sffamily\large\bfseries\color{battleshipgrey}}


\renewcommand*{\chapterpagestyle}{fancy}


\pagestyle{fancy}
\lhead{Skript Technische Informatik}
\chead{}
\rhead{\leftmark}
\lfoot{Christian Kroh \texttt{(s1428123)}}
\cfoot{}
\rfoot{\thepage}



\lstset{%
    columns=fullflexible,
    aboveskip=5pt,
    belowskip=10pt,
    basicstyle=\small\ttfamily,
    numbers=left,
    stepnumber=1, 
    numbersep=13pt,
    showspaces=false,
    showstringspaces=false,
    showtabs=false,
    xleftmargin=20pt,
    xrightmargin=10pt,
    framesep=5pt,
    framerule=3pt,
    frame=leftline,
    tabsize=2,
    breaklines=true,
    breakatwhitespace=true,
}



\title{Einführung in die Technische Informatik\\Skript}
\subtitle{In der Hoffnung, dass es was nützt ...}
\author{Christian Kroh}
\date{\today{}, Dresden}

\begin{document}
\maketitle
\newpage
\tableofcontents
\newpage

\chapter{Allgemeine Einführung}
\section{Qualifikationsziele}
Die Studierenden kennen Systemarchitekturen und Modellierungsparadigmen
von VLSI-Systemen.\\
Sie sind in der Lage Beschreibungen von Hardware-Systemen durch
Simulation zu verifizieren und mithilfe typischer Werkzeuge in reale
Schaltungen umzuwandeln.\\
Sie können den Ressourcenbedarf, das Zeitverhalten und die Verlustleistung abschätzen oder evaluieren und daraus Entwurfsentscheidungen ableiten.
\section{Literatur}
\begin{itemize}
\item F. Kesel und R. Bartholomä: Entwurf von digitalen Schaltungen und
	Systemen mit HDLs und FPGAs, Oldenbourg Wissenschaftsverlag,
	ISBN 978-3-486-58976-4.
\item H.-D. Wuttke und K. Henke: Schaltsysteme – Eine automatenorientierte
	Einfürung, Pearson Studium, ISBN 3-8273-7035-3.
\item H. M. Lipp und J. Becker: Grundlagen der Digitaltechnik,
Oldenbourg Wissenschaftsverlag, ISBN 978-3-486-59747-9.

\end{itemize}



\newpage
\part{VLSI-Systementwurf}
\newpage
\chapter{Einführung}
\section{Themenschwerpunkte}
	\begin{description}

		\item[Verarbeitungsleistung]\hfill
			\\\textbf{Im Vordergrund stehen:}
			\begin{itemize}
				\item Schnelle Verarbeitung auch einzelner Bits
				\item Parallelität auf: \hfill
				\begin{itemize}
					\item Bitebene
					\item Befehlseben
					\item Threadebene
					\item Prozess- und Anwendungsebene
				\end{itemize}
				\item dynamische Rekonfiguration
			\end{itemize}
			$\Rightarrow$ Erfüllung der gegebenen Anforderungen

		\item[Systemintegration]\hfill
			\\\textbf{Im Vordergrund stehen:}
			\begin{itemize}
				\item Mehrprozessorsyteme, Mehrkern-, Vielkernprozessorsysteme
				\item Mehr-Chip- / Einzel-Chip-Lösungen (System-on-a-Chip)
				\item Parallele Entwicklung von HW und SW
	(HW-/SW-Codesign)
				\item System-Prototyping (FPGA-Entwurf)
			\end{itemize}
			$\Rightarrow$ Kosteneinsparung, Entwicklungszeiteinsparung (Time-to-Market)

		\item[Verlustleistung]
			\hfill\\\textbf{Im Vordergrund stehen:}
			\begin{itemize}
				\item Verlustleistung im Standby (Akkubetrieb)
				\item Maximales Abwärmebudget
			\end{itemize}

			\hfill\\\textbf{Kenngrößen:}
			\begin{itemize}
				\item Statische und dynamische Verlustleistung
				\item MIPS pro Watt
			\end{itemize}

			$\Rightarrow$ Sowohl für eingebettete Systeme als auch für Server


		\item[Korrektheit]
			\hfill\\\textbf{Aspekte:}
			\begin{itemize}
				\item Verifikation eines Schaltkreises, simulativ / formal
				\item Profiling und Debugging unter Echtzeitbedingungen (Trace)
			\end{itemize}
			$\Rightarrow$ Fehlerfreier Erstentwurf


		\item[Fehlertoleranz]\hfill
			\\\textbf{Toleranz gegenüber:}
			\begin{itemize}
				\item Permanenten Fehlern (zeitunabhängig nach erstem Auftreten)
				\item Intermitierenden Fehlern (nur unter bestimmten Betriebsbedingungen)
				\item Transienten Fehlern (aufgrund statischer Störungen)
			\end{itemize}

			\hfill\\\textbf{Fehlererkennung und -korrektur:}
			\begin{itemize}
				\item Autonom durch Hardwarearchitektur
				\item Aus Kombination von HW und SW
			\end{itemize}

			$\Rightarrow$ Insbesondere wichtig für Sicherheitskritische, hochverfügbare und langlebige zuverlässige Systeme
			\\$\Rightarrow$ Steigende Siginifikanz mit abnehmenden Strukturgrößen (Integrationsgrad) sowie steigender Transistoranzahl (Moore's Law)
	\end{description}


\section{Inhalte der Lehrveranstaltung}
\subsection{Inhalte der Vorlesung}
	\begin{enumerate}
		\item Klassifikation von Schaltkreisen
		\item Grundlagen des Schaltkreisentwurfs
		\item Automatendarstellung, -kopplung, -vereinfachung
		\item Hardwarebeschreibungssprachen
		\item Programmierabare Schaltkreise, insbesondere FPGAs - Teil 1
		\item Programmierabare Schaltkreise, insbesondere FPGAs - Teil 2
		\item Modellierung und Simulation
		\item Zeitverhalten und Test
		\item Hochgeschwindigkeit und Verlsustleistung
		\item Anwendungsbeispiele
	\end{enumerate}

\subsection{Inhalte des Praktikums}
	\begin{enumerate}
		\item Altera Quartus-Toolchain \& Praktikumsboard DE0 mit Cyclone-3
		\item Schaltnetze und Schaltwerke
		\item Modularisierung
		\item \grqq Komplexe\grqq Anwendung: Stoppuhr
	\end{enumerate}


\section{Klassifikation von ICs}

\subsection{Zwei Sichten}
	\hfill\\Klassifizierung von integrierten Schaltkreisen (ICs) in
	\begin{itemize}
		\item Standarschaltkreis (Standard-IC) und
		\item applikationspezifische Schaltkreise (Application-Specific IC, \textbf{ASIC})
	\end{itemize}
	unter zwei Gesichtspunkten möglich:
	\begin{itemize}
		\item \textbf{Herstellungssicht}
		\item \textbf{Entwurfssicht}
	\end{itemize}

	\paragraph{Herstellungssicht:}
	\begin{itemize}
		\item Standard-IC = große Stückzahl für viele Kunden
		\item ASIC = für eine/n Kunden/Applikation speziell entwickelter und gefertigter IC mit zugeschnittener Funktionalität
	\end{itemize}

	\paragraph{Entwurfssicht:}
	\begin{itemize}
		\item Standard-IC = (Hardware-)Funktionalität kann nicht vom Anwender beeinflusst werden
		\item ASIC = vom Anwender selbst entwickelte (Hardware-)Funktionalität
	\end{itemize}

	\newpage
	\begin{center}
	\includegraphics[page=17, width=0.7\linewidth, trim=30mm 20mm 20mm 40mm, clip]{\Path/resources/Vorlesung/VLSI/01_Einfuehrung.pdf}
	\\\includegraphics[page=18, width=0.7\linewidth, trim=30mm 20mm 20mm 40mm, clip]{\Path/resources/Vorlesung/VLSI/01_Einfuehrung.pdf}
	\\\includegraphics[page=19, width=0.7\linewidth, trim=30mm 20mm 20mm 40mm, clip]{\Path/resources/Vorlesung/VLSI/01_Einfuehrung.pdf}
	\end{center}

\subsection{Anwenderprogrammierbare IC (ASIC)}
	\paragraph{Merkmale:}
	\begin{itemize}
		\item Field-Programmable $\Leftrightarrow$ feldprogrammierbar
		\item Vor Ort (im Feld) vom Anwender programmierbar
		\item Hardware ist fix. Funktionalität kann aber mittels spezieller Konfiguration \grqq programmiert\grqq werden
	\end{itemize}
	\paragraph{Anwendung:}
	\begin{itemize}
		\item Anwendungsspezifische IC bei kleinen und mittleren Stückzahlen
		\item Mehrfach neu programmierbar zwecks Optimierung und Fehlerbehebung, auch während des praktischen Einsatzes
		\item Einfache Integration eines ganzen Systems auf einem Chip
		\item Prototyping, HW-/SW-Codesign
	\end{itemize}

\subsection{Hardwareprogrammierung}
	\begin{center}
	\includegraphics[page=21, width=0.7\linewidth, trim=30mm 20mm 20mm 40mm, clip]{\Path/resources/Vorlesung/VLSI/01_Einfuehrung.pdf}
	\end{center}
%\newpage
\subsection{Programmiertechnologien}
	\begin{center}
	\includegraphics[page=22, width=0.7\linewidth, trim=30mm 20mm 20mm 40mm, clip]{\Path/resources/Vorlesung/VLSI/01_Einfuehrung.pdf}
	\\\includegraphics[page=23, width=0.7\linewidth, trim=30mm 20mm 20mm 40mm, clip]{\Path/resources/Vorlesung/VLSI/01_Einfuehrung.pdf}
	\\\includegraphics[page=24, width=0.7\linewidth, trim=30mm 20mm 20mm 40mm, clip]{\Path/resources/Vorlesung/VLSI/01_Einfuehrung.pdf}
	\\\includegraphics[page=25, width=0.7\linewidth, trim=30mm 20mm 20mm 40mm, clip]{\Path/resources/Vorlesung/VLSI/01_Einfuehrung.pdf}
	\\\includegraphics[page=26, width=0.7\linewidth, trim=30mm 20mm 20mm 40mm, clip]{\Path/resources/Vorlesung/VLSI/01_Einfuehrung.pdf}
	\\\includegraphics[page=27, width=0.7\linewidth, trim=30mm 20mm 20mm 40mm, clip]{\Path/resources/Vorlesung/VLSI/01_Einfuehrung.pdf}
	\\\includegraphics[page=28, width=0.7\linewidth, trim=30mm 20mm 20mm 40mm, clip]{\Path/resources/Vorlesung/VLSI/01_Einfuehrung.pdf}
	\end{center}

\subsection{Klassifikation Anwenderprogrammierbare IC}
	\begin{center}
	\includegraphics[page=29, width=0.7\linewidth, trim=30mm 20mm 20mm 40mm, clip]{\Path/resources/Vorlesung/VLSI/01_Einfuehrung.pdf}
	\\\includegraphics[page=30, width=0.7\linewidth, trim=30mm 20mm 20mm 40mm, clip]{\Path/resources/Vorlesung/VLSI/01_Einfuehrung.pdf}
	\\\includegraphics[page=31, width=0.7\linewidth, trim=30mm 20mm 20mm 40mm, clip]{\Path/resources/Vorlesung/VLSI/01_Einfuehrung.pdf}
	\end{center}

\newpage

\chapter{Schaltkreisentwurf}

\section{Abstraktionsebenen und Sichten}
	\paragraph{Ebenen des Entwurfs:}
	\begin{itemize}
		\item Charakterisierung des jeweiligen Detailiertheitsgrades der Beschreibung des Entwurfsgrades
		\item Abstraktionsgrad von der eigentlichen physikalischen Realisierung
		\item Abstraktionsniveaus, Hierarchien
	\end{itemize}

	\paragraph{Sichten des Entwurfs:}
	\begin{itemize}
		\item Betrachtung des Entwurfsgegenstandes aus verschiedenen Richtungen
		\item Sicht = Eigenschaften die den Entwurfsgegenstand Charakterisieren
		\item Alle Sichten auf jeder Entwurfseben $\Rightarrow$ Y-Diagramm / x-Diagramm
	\end{itemize}

\subsection{Abstraktionsebenen}
	\paragraph{Systemebene (system level):} Systemkonzept des Entwurfsgegenstandes
	\paragraph{Algorithmische Ebene (algorithm level):} Algorithmische Beschreibung des Entwurfsgegenstandes
	\paragraph{Register-Transfer Ebene (register-transfer-level - RTL):} Datentransfer und -verarbeitung zwischen Registern
	\paragraph{Logikebene (gate level):} Beschreibung auf Gatterniveau
	\paragraph{Schaltkreiseben:} Transistorebene im weiteren Sinne, umfasst:
		\begin{itemize}
			\item Schalterebene
			\item Schaltungsebene
			\item Bauelementebene
			\item Technologieebene
		\end{itemize}

\subsection{Sichten}
	\paragraph{Verhaltenssicht:} Beschreibung des zeitlichen Verhaltens durch charaktersierende Variablen und deren Werteverläufe über die Zeit
		
		\begin{center}
			\boldmath\( \mathlarger{\vec{y}(t) = f(\vec{x}(t))} \)
		\end{center}
	\paragraph{Struktursicht:} Spezifizierung eines Objektes durch Subobjekte und deren Verbindungsstrukturen
	
	\paragraph{Geometriesicht:} Räumliche Ausdehnung der Subobjekte
	
	\paragraph{Testsicht:} Existenz oder Nichtexistenz angenommener struktureller oder funktionaler Defekte \\(F. J. Rammig, Systematischer Entwurf digitaler Systeme, B.G Teubner Stuttgart 1989)
	
\subsection{Y-Diagramm nach Gajski}
	\begin{itemize}
		\item Stellt Ebenen und Sichten in FOrm eines Y-Diagrammes dar
		\item Ursprünglich nur 3 Ebenen
		\item Heute: Erweiterung auf 5 Ebenen:
			\begin{itemize}
				\item Systemebene
				\item Algorithmische Ebene
				\item Register-Transfer Ebene
				\item Logikeben
				\item Schaltkreisebene
			\end{itemize}
	\end{itemize}
	\begin{center}
		\includegraphics[page=9, width=0.8\linewidth, trim=30mm 20mm 20mm 40mm, clip]{\Path/resources/Vorlesung/VLSI/02_Schaltkreisentwurf.pdf}
		\\\includegraphics[page=10, width=0.8\linewidth, trim=30mm 20mm 20mm 40mm, clip]{\Path/resources/Vorlesung/VLSI/02_Schaltkreisentwurf.pdf}
	\end{center}
	
\newpage
\section{Entwurfsablauf}
	\paragraph{Allgemein:} Transformation einer Aufgabenstellung (Pflichtenheft) in einen fertigen Schaltkreis
	\paragraph{Top-Down-Strategie:}
		\begin{itemize}
			\item Systemebene $\rightarrow$ Schaltkreisebene
			\item Vorteil: Parallele Entwicklung auf unteren Ebenen
			\item Nachteil: Systemspezifikation zu Projektbeginn oft zu ungenau
		\end{itemize}
	\paragraph{Bottom-Up-Strategie:}
		\begin{itemize}
			\item Analyse vorhandener Komponenten
			\item Zusammensetzen von neuen Komponenten auf höherer Ebene im Sinne der Aufgabenstellung
			\item Nachteil: globales Ziel wird nicht immer erreicht
		\end{itemize}
	\paragraph{$\Rightarrow$ Meet-in-the-Middle}
	
	\paragraph{Entwurfsschritt:}
		\begin{itemize}
			\item generierende Aktivität
			\item überprüfende Aktivität
		\end{itemize}
	\paragraph{Syntheseschritt:}
		\begin{itemize}
			\item Abbildung eines Entwurfsschrittes in Richtung auf das Entwurfsziel
			\item Abstraktionsgrad sinkt, Detailhiertheitsgrad steigt
			\item Einbringung neuer Informationen
		\end{itemize}
	\paragraph{Analyseschritt:}
		\begin{itemize}
			\item Abbildung eines Entwurfsschrittes in umgekehrter Richtung zum Syntheseschritt
			\item Gewinnung abstrakter Informationen durch Zusammenfassen und Generalisieren von Details (Extraktionsprozess)
			\item Beispiel: Validierung eines Syntheseschrittes
		\end{itemize}
		
	\begin{center}
		\includegraphics[page=13, width=0.39\linewidth, trim=40mm 30mm 40mm 50mm, clip]{\Path/resources/Vorlesung/VLSI/02_Schaltkreisentwurf.pdf}
		\includegraphics[page=14, width=0.6\linewidth, trim=0mm 30mm 0mm 45mm, clip]{\Path/resources/Vorlesung/VLSI/02_Schaltkreisentwurf.pdf}
	\end{center}
	
\newpage
\section{Entwurfsstile}
	\begin{center}
		\includegraphics[page=15, width=0.8\linewidth, trim=40mm 28mm 40mm 60mm, clip]{\Path/resources/Vorlesung/VLSI/02_Schaltkreisentwurf.pdf}
	\end{center}
	
	\subsubsection{Gegenüberstellung Entwurfsalternativen}
		\begin{center}
			\includegraphics[page=16, width=0.8\linewidth, trim=35mm 28mm 30mm 60mm, clip]{\Path/resources/Vorlesung/VLSI/02_Schaltkreisentwurf.pdf}
		\end{center}

\subsection{Full-Custom Entwurf}
	\begin{center}
		\includegraphics[page=17, width=0.8\linewidth, trim=40mm 28mm 40mm 60mm, clip]{\Path/resources/Vorlesung/VLSI/02_Schaltkreisentwurf.pdf}
	\end{center}
	\paragraph{Merkmale:}
	\begin{itemize}
		\item Platzierung und Verdrahtung selbst entworfener Transistoren \& Gatter
		\item Auch Mischung von Analog- und Digitaltechnik, z.B. für spezielle I/O-Signaltreiber
		\item Erfüllung spezieller Anforderungen, z.B. gehärtet gegen Strahlung
		\item Häufig nur auf kleine Teilschaltungen angewendet
		\item Hoch qualifizierte Entwicklungsingenieure mit Detailkenntnissen zu den Prozessen erforderlich
	\end{itemize}
	\paragraph{Anwendung:}
	\begin{itemize}
		\item Sensorik, Mixed-Signal-Schaltungen
		\item Raumfahrt
	\end{itemize}
	\subsubsection{Beispiel: Inverter}
		\begin{center}
			\includegraphics[page=19, width=0.8\linewidth, trim=40mm 28mm 40mm 60mm, clip]{\Path/resources/Vorlesung/VLSI/02_Schaltkreisentwurf.pdf}
		\end{center}
	
\newpage
\subsection{Standardzellenentwurf}
	\begin{center}
		\includegraphics[page=20, width=0.8\linewidth, trim=40mm 28mm 40mm 60mm, clip]{\Path/resources/Vorlesung/VLSI/02_Schaltkreisentwurf.pdf}
	\end{center}
	\paragraph{Merkmale:}
	\begin{itemize}
		\item Platzierung und Verdrahtung vorgegebener Gatter- oder Makrozellen
		\item Auswahl einer Technologiebibliothek nach:
			\begin{itemize}
				\item Fertigungstechnologie, Strukturbreite und Funktion
				\item High-Speed, Low-Leakage oder Mischung aus Beidem
			\end{itemize}
		\item Werkzeuggestützte Platzierung und Verdrahtung
		\item Makrozellen:
			\begin{itemize}
				\item Vordefiniert oder per Generator kundenspezifisch erzeugt
				\item Bsp: RAM, ALU, I/O-Komponenten
			\end{itemize}
	\end{itemize}
	\paragraph{Anwendung:}
	\begin{itemize}
		\item Anwendungsspezifische Schaltkreise mit hohen Stückzahlen
		\item Starke Optimierung bzgl. Chipfläche, Geschwindigkeit und Verlustleistung
	\end{itemize}
	\subsubsection{Konventionelle Architektur vs. Strukturierte Architektur}
		\begin{center}
			\includegraphics[page=22, width=0.45\linewidth, trim=60mm 28mm 60mm 60mm, clip]{\Path/resources/Vorlesung/VLSI/02_Schaltkreisentwurf.pdf}
			\includegraphics[page=23, width=0.45\linewidth, trim=60mm 28mm 60mm 60mm, clip]{\Path/resources/Vorlesung/VLSI/02_Schaltkreisentwurf.pdf}
		\end{center}
	
\newpage
\subsection{Maskenprogrammierbare Gate Arrays}
	\begin{center}
		\includegraphics[page=27, width=0.8\linewidth, trim=40mm 28mm 40mm 60mm, clip]{\Path/resources/Vorlesung/VLSI/02_Schaltkreisentwurf.pdf}
	\end{center}
	\paragraph{Merkmale:}
	\begin{itemize}
		\item Festlegung der Funktion mittels Masken bei der Halbleiterfertigung
		\item Ausgewählte Masken: teilweise vorgefertigte IC (Master)
		\item Beispiele:
			\begin{itemize}
				\item MPGA (Maskprogrammable Gate-Array): auch MGA
					\begin{itemize}
						\item Vorgefertigte universelle Gatter/Makros mit fester Anordnung
						\item Verdrahtung kundenspezifisch
					\end{itemize}
				\item MROM: ROM dessen Inhalt vom Kunden mit Masken festgelegt wird
			\end{itemize}
	\end{itemize}
	\paragraph{Anwendung:}
	\begin{itemize}
		\item Anwendungsspezifische Schaltkreise bei mittleren Stückzahlen
		\item Kostenoptimiert mit reduzierten Optimierungspotential
	\end{itemize}
	\subsubsection{MPGA}
		\begin{center}
			\includegraphics[page=29, width=0.45\linewidth, trim=60mm 28mm 60mm 60mm, clip]{\Path/resources/Vorlesung/VLSI/02_Schaltkreisentwurf.pdf}
		\end{center}
	
	
\newpage
\subsection{Anwenderprogrammierbare IC}
	\begin{center}
		\includegraphics[page=30, width=0.8\linewidth, trim=40mm 28mm 40mm 60mm, clip]{\Path/resources/Vorlesung/VLSI/02_Schaltkreisentwurf.pdf}
	\end{center}
	\paragraph{Merkmale:}
	\begin{itemize}		
		\item Field-Programmable $\Leftrightarrow$ feldprogrammierbar
		\item Vor Ort (im Feld) vom Anwender programmierbar
		\item Hardware ist fix. Funktionalität kann aber mittels spezieller Konfiguration \grqq programmiert\grqq werden
	\end{itemize}
	\paragraph{Anwendung:}
	\begin{itemize}
		\item Anwendungsspezifische IC bei kleinen und mittleren Stückzahlen
		\item Mehrfach neu programmierbar zwecks Optimierung und Fehlerbehebung, auch während des praktischen Einsatzes
		\item Einfache Integration eines ganzen Systems auf einem Chip
		\item Prototyping, HW-/SW-Codesign
	\end{itemize}
	\subsubsection{Hardwareprogrammierung}
		\begin{center}
		\includegraphics[page=21, width=0.7\linewidth, trim=30mm 20mm 20mm 40mm, clip]{\Path/resources/Vorlesung/VLSI/01_Einfuehrung.pdf}
		\end{center}
	\subsubsection{Klassifikation Anwenderprogrammierbare IC}
		\begin{center}
			\includegraphics[page=33, width=0.8\linewidth, trim=35mm 28mm 30mm 58mm, clip]{\Path/resources/Vorlesung/VLSI/02_Schaltkreisentwurf.pdf}
		\end{center}
	\subsubsection{CPLD}
		Globale Vernetzung einer kleiner Anzahl von Funktionsblöcken (FB)
		\begin{center}
			\includegraphics[page=34, width=0.8\linewidth, trim=35mm 23mm 30mm 70mm, clip]{\Path/resources/Vorlesung/VLSI/02_Schaltkreisentwurf.pdf}
		\end{center}
	\subsubsection{CPLD-Funktionsblock}
		Funktionsblock bestehend aus PLA (Und/Oder-Matrix) und Makrozellen (MC)
		\begin{center}
			\includegraphics[page=35, width=0.5\linewidth, trim=35mm 23mm 30mm 70mm, clip]{\Path/resources/Vorlesung/VLSI/02_Schaltkreisentwurf.pdf}
		\end{center}
	\subsubsection{PLA}
		\begin{center}
			\includegraphics[page=36, width=0.8\linewidth, trim=35mm 23mm 30mm 40mm, clip]{\Path/resources/Vorlesung/VLSI/02_Schaltkreisentwurf.pdf}
		\end{center}
	\subsubsection{Makrozellen + I/O-Block}
		Verschiedene Betriebsspannungen für digitale Lokig (Core) und I/O-Pads
		\begin{center}
			\includegraphics[page=37, width=0.8\linewidth, trim=35mm 23mm 30mm 70mm, clip]{\Path/resources/Vorlesung/VLSI/02_Schaltkreisentwurf.pdf}
		\end{center}
		\paragraph{Konfiguration der Makrozelle:}	
		\begin{center}
			\includegraphics[page=38, width=0.8\linewidth, trim=35mm 93mm 30mm 70mm, clip]{\Path/resources/Vorlesung/VLSI/02_Schaltkreisentwurf.pdf}
		\end{center}
		\paragraph{Steuerung des I/O-Blocks:}\hfill\\
		Umschaltung des I/O-Pins zwischen Ein- und Ausgang (Tri-State) zur Laufzeit mittels seperatem Produktterm möglich.
	
	\subsubsection{FPGA-Architektur}	
		\begin{center}
			\includegraphics[page=39, width=0.8\linewidth, trim=35mm 30mm 20mm 58mm, clip]{\Path/resources/Vorlesung/VLSI/02_Schaltkreisentwurf.pdf}
		\end{center}
	
\section{Entwurfswerkzeuge}
	\begin{center}
		\includegraphics[page=40, width=0.8\linewidth, trim=35mm 87mm 30mm 65mm, clip]{\Path/resources/Vorlesung/VLSI/02_Schaltkreisentwurf.pdf}
	\end{center}
	Auswahl CAD-Werkzeuge:
	\begin{itemize}
		\item Cadence
		\item Xilinx ISE
		\item Altera Quartus
		\item Synopsys
	\end{itemize}
	
	
	
	
	

\chapter{Automaten}

\section{Automatendarstellung}

\subsection{Betrachtungsweisen}
	\paragraph{Verschiedene Sichten/Semantik:}
	\begin{itemize}
		\item Zustandsübergangsdiagramm (state diagram)
			\begin{itemize}
				\item Vernetzung von Zuständen
				\item Beispiele: UML-Zustandsdiagramm, Automatengrpahen, SM Charts, GRAFCET, Sequential Function Charts (SFC)
			\end{itemize}
		\item Ablaufdiagramm (flow chart)
			\begin{itemize}
				\item Vernetzung von Prozessen
				\item Zustand ergibt sich aus der Verkettung aller Variablenzustände
				\item Beispiele: UML-Aktivitätsdiagramm, Programmablaufplan (PAP)
			\end{itemize}
	\end{itemize}
	
	\subsubsection{Binärzähler mod 4 mit Trigger X und Übertrag Y}	
		\begin{center}
			\includegraphics[page=4, width=0.6\linewidth, trim=35mm 27mm 30mm 58mm, clip]{\Path/resources/Vorlesung/VLSI/03_Automaten.pdf}
		\end{center}
	
\subsection{Automatengraphen}	
	\begin{center}
		\includegraphics[page=5, width=0.8\linewidth, trim=35mm 67mm 30mm 58mm, clip]{\Path/resources/Vorlesung/VLSI/03_Automaten.pdf}
	\end{center}
	\paragraph{Endlicher Automat:} Menge der möglichen Eingabezeichen, Ausgabezeichen und inneren Zustände ist endlich ...
	\paragraph{Synchron getakteter Automat:} Zustandsübergänge aller Speicherglieder erfolgen gleichzeitig, synchron zu einem Taktsignal
	
	\subsubsection{Notation}
		\begin{center}
			\includegraphics[page=6, width=0.65\linewidth, trim=35mm 37mm 30mm 58mm, clip]{\Path/resources/Vorlesung/VLSI/03_Automaten.pdf}
		\end{center}
	\subsubsection{Grafische Darstellung}
		\begin{center}
			\includegraphics[page=7, width=0.65\linewidth, trim=35mm 39mm 30mm 58mm, clip]{\Path/resources/Vorlesung/VLSI/03_Automaten.pdf}
		\end{center}
		\paragraph{$\Rightarrow$ Prüfung auf:} Vollständigkeit und Widerspruchfreiheit
	
	\subsubsection{Getaktete Automaten}
		\paragraph{Theoretische Informatik}
		\begin{itemize}
			\item Verarbeitung von EIngabezeichen sofern vorhanden
			\item Jedes Zeichen in der Eingabe wird einmalig verarbeitet
		\end{itemize} 
		\paragraph{Technische Informatik}\hfill\\
		Taktung des Automaten mit einem Taktsignal $\rightarrow$ Abtatstung der Eingabe\\
		Konsequenzen:
		\begin{itemize}
			\item Zeichen für \grqq keine Eingabe \grqq erforderlich
			\item Abtastung \grqq derselben Eingabe \grqq in aufeinanderfolgenden Takten möglich
		\end{itemize} 
		$\Rightarrow$ Korrekte Modellierung des Taktsignals erforderlich
	
		\paragraph{Beispiel - Taktzustandsgesteuertes D-Flip-Flop (Latch)}\hfill\\
		(Taktsignal C als Eingabe) 
		\begin{center}
			\includegraphics[page=9, width=0.65\linewidth, trim=35mm 50mm 30mm 80mm, clip]{\Path/resources/Vorlesung/VLSI/03_Automaten.pdf}
		\end{center}
	
		\paragraph{Beispiel - Taktflankengesteuertes D-Flip-Flop}\hfill\\
		(Taktsignal C als Eingabe) 
		\begin{center}
			\includegraphics[page=10, width=0.65\linewidth, trim=35mm 26mm 30mm 80mm, clip]{\Path/resources/Vorlesung/VLSI/03_Automaten.pdf}
		\end{center}
	
		\paragraph{Beispiel - Taktflankengesteuertes D-Flip-Flop}\hfill\\
		(Definition: Zustandsübergang nur bei taktflanke) 
		\begin{center}
			\includegraphics[page=11, width=0.65\linewidth, trim=35mm 50mm 30mm 80mm, clip]{\Path/resources/Vorlesung/VLSI/03_Automaten.pdf}
		\end{center}
	
\subsection{SM-Charts}
	\begin{center}
		\includegraphics[page=12, width=0.65\linewidth, trim=35mm 30mm 30mm 60mm, clip]{\Path/resources/Vorlesung/VLSI/03_Automaten.pdf}
	\end{center}
	
	\subsubsection{Verzweigungen}
		\paragraph{Weitere Merkmale:}
		\begin{itemize}
			\item Selbstdefinierte Abkürzungen für komplexe Ausdrücke üblich
			\item im Allgemeinen nur Prüfung einer Eingangsvariablen (nicht Eingabezeichen) pro Verzweigung
		\end{itemize}
		\paragraph{Weitere Beispiele:}
		\begin{center}
			\includegraphics[page=13, width=0.65\linewidth, trim=35mm 27mm 30mm 111mm, clip]{\Path/resources/Vorlesung/VLSI/03_Automaten.pdf}
		\end{center}
		
	
	\subsubsection{Ausgaben}
		\paragraph{Ebenso:}
		Selbstdefinierte Abkürzungen für komplexe Ausdrücke üblich
		
		\paragraph{Weitere Beispiele:}
		\begin{center}
			\includegraphics[page=14, width=0.65\linewidth, trim=35mm 27mm 30mm 110mm, clip]{\Path/resources/Vorlesung/VLSI/03_Automaten.pdf}
		\end{center}
	
	\subsubsection{Vorteile gegenüber Automatengraphen}
		\begin{itemize}
			\item prinzipiell Vollständig
			\item präzise Modellierung hierarchischer Verzweigungen
		\end{itemize}
		\begin{center}
			\includegraphics[page=15, width=0.65\linewidth, trim=35mm 27mm 30mm 80mm, clip]{\Path/resources/Vorlesung/VLSI/03_Automaten.pdf}
		\end{center}
			

\subsection{GRAFCET \& SFC}
	\paragraph{Allgemeines:}
	\begin{itemize}
		\item Struktur entspricht einem Petri-Netz
		\item Darstellung von schrittweise ausgeführten Ablaufbeschreibungen
		\item Anwendung in Automatisierungstechnik und Verfahrenstechnik
		\item Verwendbar zur Programmierung von Speicherprogrammierbaren Steuerungen (SPS)
		\item Automatenkopplung nicht vorgesehen; muss durch Eingangs- und Ausgangsvariablen realisiert werden
	\end{itemize}
	\hfill\\
	\begin{tabular}{ll}
		GRAFCET & GRAphe Functionnel de Commande Etapes/Transitions\\
		SFC & sequential function chart
	\end{tabular}
	
	
	\subsubsection{Grundelemente}
		\begin{itemize}
			\item Schritt:
				\begin{itemize}
					\item entspricht einem Zustand
					\item mindestens ein Initialisierungsschritt notwendig
				\end{itemize}
			\item Transition
				\begin{itemize}
					\item Schaltbedingung für Übergang zwischen zwei Schritten (boolesche Gleichung) = zeitliche Ereignisse
					\item Schritte und Transitionen folgen immer aufeinander
					\item Leserichtung: Transitionen können nur von oben nach unten durchlaufen werden
				\end{itemize}
			\item Aktion
				\begin{itemize}
					\item einem Schritt zugeordnet
					\item realisiert die Ausgabe
					\item verschiedene Aktionsarten möglich
				\end{itemize}
		\end{itemize}
	
	\subsubsection{Aktionen}
		\begin{center}
			\includegraphics[page=18, width=0.65\linewidth, trim=85mm 27mm 50mm 43mm, clip]{\Path/resources/Vorlesung/VLSI/03_Automaten.pdf}
		\end{center}
		
	\subsubsection{Kontrollfluss: Verzweigungen}
		\begin{center}
			\includegraphics[page=19, width=0.65\linewidth, trim=55mm 27mm 50mm 60mm, clip]{\Path/resources/Vorlesung/VLSI/03_Automaten.pdf}
		\end{center}
		
	\subsubsection{Kontrollfluss: Schleife}
		\begin{center}
			\includegraphics[page=20, width=0.65\linewidth, trim=45mm 27mm 50mm 60mm, clip]{\Path/resources/Vorlesung/VLSI/03_Automaten.pdf}
		\end{center}
		
	\subsubsection{Vergleich von GRAFCET \& SFC}
		\begin{itemize}
			\item identische Grundstruktur, Verzweigung und Initialisierung
			\item Ablauf-Ends in GRAFCET nicht notwendig, dann SPrung zu Initialisierungsschritt
			\item Schleifen in SFC
				\begin{itemize}
					\item Vorwärtssprung = Alternativverzweigung
					\item Rückwärtssprung = Schleife
				\end{itemize}
			\item Aktionen s.o.
		\end{itemize}
		$\Rightarrow$ GRAFCET ist für den Entwurf konzipiert, SFC für die Implementierung
		
	\subsubsection{Beispiel zu GRAFCET}
		\begin{center}
			\includegraphics[page=22, width=0.65\linewidth, trim=100mm 27mm 50mm 42mm, clip]{\Path/resources/Vorlesung/VLSI/03_Automaten.pdf}
		\end{center}
		


\section{Automatenkopplung}
	\paragraph{Zerlegung der Gesamtaufgabe in Teilschaltungen zwecks:}
	\begin{itemize}
		\item Nebenläufigkeit\\(statt sequentieller Abarbeitung)
		\item Energieeinsparung ducrh Abschaltung ungenutzter Komponenten\\(statt Dauerbetrieb der gesamten Schaltung)
		\item komponentenbasierter, testfreundlicher Entwurf\\(statt monolitischem Design)
		\item Integration von Komponenten von Drittanbietern\\(statt eigenem Entwurf aller Komponenten)
	\end{itemize}
	
	\paragraph{Außerdem:} 
	Kommunikation mit der Außenwelt erfordert im Allgemeinen nebenläufig arbeitende I/O-Controller aufgrund asynchron eintreffender Ereignisse (Nachrichten) $\rightarrow$ Client-Server Architektur
	
	\subsubsection{Zeitliche Kopplung}
		\begin{center}
			\includegraphics[page=25, width=0.65\linewidth, trim=30mm 50mm 30mm 65mm, clip]{\Path/resources/Vorlesung/VLSI/03_Automaten.pdf}
		\end{center}
		
	\subsubsection{Asynchrone Kopplung synchroner Teilschaltungen}
		\paragraph{Merkmal:} Mehrere (lokale) Taktsignale / Taktdomänen (clock domain)
		\paragraph{Vorteile:}
		\begin{itemize}
			\item Timing-Analyse sichert Zeitverhalten innerhalb der Domäne
			\item passende Taktfrequenz innerhalb der Domäne
			\begin{itemize}
				\item Kombination High-Speed und Low-Power
				\item Taktfrequenz dynamisch anpassbar
				\item Abschalten einer Domäne für Standby
			\end{itemize}
		\end{itemize}
		\paragraph{Nachteil:}\hfill\\
		Datenaustausch zwischen Taktdomänen erfordern:
		\begin{itemize}
			\item Module mit Single-Bit-Synchronizer, Cross-Clock-FIFOs
			\item spezielle Timing-Constraints
		\end{itemize}
	
	\subsubsection{Synchrone Kopplung mit mehreren Taktsignalen}
		\paragraph{Merkmal:} Mehrere (lokale) Taktsignale / Taktdomänen, zwischen denen aber spezielle Abhängigkeiten bestehen (dependent clocks)
		\paragraph{Vorteile:}
		\begin{itemize}
			\item Timing-Analyse sichert Zeitverhalten innerhalb der Domäne und auch zwischen den Domänen
			\item passende feste Taktfrequenz je Domäne
			\item keine speziellen Synchronisationselemente erforderlich
		\end{itemize}
		\paragraph{Nachteil:} Taktfrequenzen sind statisch.
		
		
\subsection{Synchrone Kopplung}
	Kommunikation über:
	\begin{itemize}
		\item Zustände oder
		\item Ausgaben
	\end{itemize}
	
	Anordnung:
	\begin{itemize}
		\item parallel oder
		\item seriell
	\end{itemize}
	
	
	\subsubsection{Kommunikation über Zustände}
		\begin{center}
			\includegraphics[page=30, width=0.65\linewidth, trim=30mm 27mm 30mm 60mm, clip]{\Path/resources/Vorlesung/VLSI/03_Automaten.pdf}
		\end{center}
	
	\subsubsection{Kommunikation über Ausgaben}
		\begin{center}
			\includegraphics[page=31, width=0.65\linewidth, trim=30mm 27mm 27mm 60mm, clip]{\Path/resources/Vorlesung/VLSI/03_Automaten.pdf}
			\includegraphics[page=32, width=0.65\linewidth, trim=30mm 27mm 30mm 60mm, clip]{\Path/resources/Vorlesung/VLSI/03_Automaten.pdf}
		\end{center}
	
	\subsubsection{Parallele Anordnung}
		\begin{center}
			\includegraphics[page=33, width=0.65\linewidth, trim=30mm 27mm 30mm 60mm, clip]{\Path/resources/Vorlesung/VLSI/03_Automaten.pdf}
		\end{center}
		
	\subsubsection{Serielle Anordnung}
		\begin{center}
			\includegraphics[page=34, width=0.65\linewidth, trim=30mm 27mm 27mm 60mm, clip]{\Path/resources/Vorlesung/VLSI/03_Automaten.pdf}
		\end{center}
		
	\subsubsection{Flusskontrolle}
		\begin{center}
			\includegraphics[page=35, width=0.65\linewidth, trim=30mm 40mm 27mm 60mm, clip]{\Path/resources/Vorlesung/VLSI/03_Automaten.pdf}
		\end{center}
		
	\subsubsection{Beispiel: Produktautomat}
		\begin{center}
			\includegraphics[page=36, width=0.65\linewidth, trim=30mm 50mm 27mm 60mm, clip]{\Path/resources/Vorlesung/VLSI/03_Automaten.pdf}
		\end{center}
		
	\subsubsection{Beispiel: Datenverarbeitung}
		\begin{center}
			\includegraphics[page=37, width=0.65\linewidth, trim=30mm 30mm 27mm 60mm, clip]{\Path/resources/Vorlesung/VLSI/03_Automaten.pdf}
			\includegraphics[page=38, width=0.65\linewidth, trim=30mm 30mm 27mm 60mm, clip]{\Path/resources/Vorlesung/VLSI/03_Automaten.pdf}
			\includegraphics[page=39, width=0.65\linewidth, trim=30mm 30mm 27mm 60mm, clip]{\Path/resources/Vorlesung/VLSI/03_Automaten.pdf}
		\end{center}
		
	\subsubsection{Beispiel: De-/Multiplexer}
		\begin{center}
			\includegraphics[page=40, width=0.65\linewidth, trim=30mm 30mm 27mm 60mm, clip]{\Path/resources/Vorlesung/VLSI/03_Automaten.pdf}
		\end{center}

\subsection{Asynchrone Kopplung}

	\subsubsection{Abtastproblem}
	
	\begin{wrapfigure}{r}{0.25\linewidth}
  		\vspace{-20pt}
	  	\begin{center}
			\includegraphics[page=41, width=0.6\linewidth, trim=210mm 45mm 27mm 65mm, clip]{\Path/resources/Vorlesung/VLSI/03_Automaten.pdf}
	  	\end{center}
		\vspace{-20pt}
	\end{wrapfigure}
	
	Kommunikation zwischen verschiednen Taktdomänen aufgrund verschiedener Taktfrequenzen und -phasen komplexer:
	\begin{itemize}
		\item keine feste Zuordnung zwischen zwei Zeitpunkten, \\allg. Annahme \( t_n = n * T_1 \neq m* T_2 = t_m\) mit $n,m \in N$
		\item Abtastung von Signalvektoren kann zu Fehlern aufgrund verschiedener Signallaufzeiten führen
	\end{itemize}
	
	\subsubsection{Flusskontrolle}
	Berücksichtigung verschiedener Taktfrequenzen im Protokoll notwendig:
	\begin{itemize}
		\item Mehrfaches Lesen desseleben Ausgabezeichens, \\wenn $f_2 > f_1$
		\item Verpassen von Ausgabezeichen, \\wenn $f_2 < f_1$
	\end{itemize} 
	
	\subsubsection{Übertragung eines einzelnen Bits}
		\begin{center}
			\includegraphics[page=43, width=0.65\linewidth, trim=30mm 27mm 27mm 60mm, clip]{\Path/resources/Vorlesung/VLSI/03_Automaten.pdf}
		\end{center}
	
	\subsubsection{Übertragung eines Signalvektors}
	Mehrere Varainten möglich:
	\begin{enumerate}
		\item Gray-Code
		\item Steuerung mittels einzelner, seperater Bits
		\item Warteschlange mittels Cross-Clock-FIFO
	\end{enumerate}
	
	\subsubsection{Variante 1: Gray-Code}
		\includegraphics[page=45, width=0.65\linewidth, trim=30mm 27mm 27mm 60mm, clip]{\Path/resources/Vorlesung/VLSI/03_Automaten.pdf}\\
		\includegraphics[page=46, width=0.65\linewidth, trim=30mm 27mm 27mm 60mm, clip]{\Path/resources/Vorlesung/VLSI/03_Automaten.pdf}
		
	\subsubsection{Variante 2: Steuerung mittels einzelner Bits}
		\includegraphics[page=47, width=0.65\linewidth, trim=25mm 30mm 27mm 60mm, clip]{\Path/resources/Vorlesung/VLSI/03_Automaten.pdf}
		
	\subsubsection{Variante 3: FIFO-Warteschlange}
		\includegraphics[page=48, width=0.65\linewidth, trim=25mm 30mm 27mm 60mm, clip]{\Path/resources/Vorlesung/VLSI/03_Automaten.pdf}\\
		
		\includegraphics[page=49, width=0.65\linewidth, trim=25mm 30mm 27mm 60mm, clip]{\Path/resources/Vorlesung/VLSI/03_Automaten.pdf}
		
		
		
		
		
\section{Initialisierung}
	
\subsection{Reset vs. Power-Up}
	\paragraph{Nicht programmierbare Schaltkreise:}
	\begin{itemize}
		\item Reset notwendig für Initialisierung der Zustandsregister
		\item Datenregister können von Automaten initialisiert werden
	\end{itemize}
	
	\paragraph{Programmierbare Schaltkreise:}
	\begin{itemize}
		\item Initiale Registerbelegung wird durch Programmierung festgelegt
		\item Reset-Eingang ist dh. optional
	\end{itemize}
	
	\paragraph{$\rightarrow$ Wiederverwendungsgerechter Entwurf:}
	\begin{itemize}
		\item Reset-Eingang vorsehen
		\item (Zustands-)Register bei Power-Up und Reset gleichermaßen belegen
	\end{itemize}
	
\subsection{Synchrones Reset}
	\paragraph{Vorteile:}
	\begin{itemize}
		\item Synthese immer möglich: Set/Reset als zusätzliche Variable in FF-Ansteuergleichung, sofern nicht explizit vorhanden
		\item Logik-Zusammenfassung möglich
		\item Überprüfung in Timing-Analys
	\end{itemize}
	
	\paragraph{Nachteil:}\hfill\\
	Free-Running Clock benötigt, damit Reset auch ausgelöst wird
	
	\paragraph{Reset-Synchronizer}
	\begin{itemize}
		\item Reset-Pin muss synchronisiert werden, damit alle Register in der gleichen Taktperiode zurückgesetzt werden
		\item Synchronisation mit (min.) 2 FFs, um Metastabilitäten zu vermeiden
		\item Reset-Pin ist ggf. zu negieren
		\item Schaltung allg. verwendbar für Synchronisation asynchroner Signale
	\end{itemize}
	\begin{center}
		\includegraphics[page=52, width=0.65\linewidth, trim=30mm 27mm 27mm 120mm, clip]{\Path/resources/Vorlesung/VLSI/03_Automaten.pdf}
	\end{center}
	
\subsection{Asynchrones Reset}
	\paragraph{Vorteile:}
	\begin{itemize}
		\item Keine Free-Running CLock benötigt um Reset auszulösen, u.U. notwendig für Register, die Schaltungsausgänge treiben
		\item Seperater FF-Eingang, damit kein Einfluss auf Timing
	\end{itemize}
	
	\paragraph{Nachteile:}
	\begin{itemize}
		\item Implementierung erfordert globale Verdrahtungsressourcen wie Taktsignale
		\item Timing-Analyse problematisch und häufig standardmäßig ignoriert
		\item Reset am Register darf nicht in zeitlicher Nähe zur Taktflanke losgelassen werden $\rightarrow$ sonst Metastabilitäten
	\end{itemize}
	
	\paragraph{Reset-Synchronizer}
	\begin{itemize}
		\item Reset muss synchron losgelassen werden, damit alle Register in der gleichen Taktperiode wieder in Betrieb gehen
		\item Reset-Pin ist ggf. zu negieren
		\item Reset-Tree sorgt für zusätzliche, notwendige! Verzögerung
	\end{itemize}
	\begin{center}
		\includegraphics[page=54, width=0.65\linewidth, trim=30mm 27mm 27mm 120mm, clip]{\Path/resources/Vorlesung/VLSI/03_Automaten.pdf}
	\end{center}
	
\subsection{Coding Guidelines für Reset}
	\paragraph{Empfehlungen:}
	\begin{itemize}
		\item Synchrones Reset
		\item Reset-Eingang setzt nur Zustandsregister zurück\\
			$\to$ Fan-Out von Reset-EIngang klein
		\item Power-Up-Wert automatisch ermitteln lassen. damit identisch zu Reset-Wert
	\end{itemize}

%\chapter{Einführung}
\section{Themenschwerpunkte}
	\begin{description}

		\item[Verarbeitungsleistung]\hfill
			\\\textbf{Im Vordergrund stehen:}
			\begin{itemize}
				\item Schnelle Verarbeitung auch einzelner Bits
				\item Parallelität auf: \hfill
				\begin{itemize}
					\item Bitebene
					\item Befehlseben
					\item Threadebene
					\item Prozess- und Anwendungsebene
				\end{itemize}
				\item dynamische Rekonfiguration
			\end{itemize}
			$\Rightarrow$ Erfüllung der gegebenen Anforderungen

		\item[Systemintegration]\hfill
			\\\textbf{Im Vordergrund stehen:}
			\begin{itemize}
				\item Mehrprozessorsyteme, Mehrkern-, Vielkernprozessorsysteme
				\item Mehr-Chip- / Einzel-Chip-Lösungen (System-on-a-Chip)
				\item Parallele Entwicklung von HW und SW
	(HW-/SW-Codesign)
				\item System-Prototyping (FPGA-Entwurf)
			\end{itemize}
			$\Rightarrow$ Kosteneinsparung, Entwicklungszeiteinsparung (Time-to-Market)

		\item[Verlustleistung]
			\hfill\\\textbf{Im Vordergrund stehen:}
			\begin{itemize}
				\item Verlustleistung im Standby (Akkubetrieb)
				\item Maximales Abwärmebudget
			\end{itemize}

			\hfill\\\textbf{Kenngrößen:}
			\begin{itemize}
				\item Statische und dynamische Verlustleistung
				\item MIPS pro Watt
			\end{itemize}

			$\Rightarrow$ Sowohl für eingebettete Systeme als auch für Server


		\item[Korrektheit]
			\hfill\\\textbf{Aspekte:}
			\begin{itemize}
				\item Verifikation eines Schaltkreises, simulativ / formal
				\item Profiling und Debugging unter Echtzeitbedingungen (Trace)
			\end{itemize}
			$\Rightarrow$ Fehlerfreier Erstentwurf


		\item[Fehlertoleranz]\hfill
			\\\textbf{Toleranz gegenüber:}
			\begin{itemize}
				\item Permanenten Fehlern (zeitunabhängig nach erstem Auftreten)
				\item Intermitierenden Fehlern (nur unter bestimmten Betriebsbedingungen)
				\item Transienten Fehlern (aufgrund statischer Störungen)
			\end{itemize}

			\hfill\\\textbf{Fehlererkennung und -korrektur:}
			\begin{itemize}
				\item Autonom durch Hardwarearchitektur
				\item Aus Kombination von HW und SW
			\end{itemize}

			$\Rightarrow$ Insbesondere wichtig für Sicherheitskritische, hochverfügbare und langlebige zuverlässige Systeme
			\\$\Rightarrow$ Steigende Siginifikanz mit abnehmenden Strukturgrößen (Integrationsgrad) sowie steigender Transistoranzahl (Moore's Law)
	\end{description}


\section{Inhalte der Lehrveranstaltung}
\subsection{Inhalte der Vorlesung}
	\begin{enumerate}
		\item Klassifikation von Schaltkreisen
		\item Grundlagen des Schaltkreisentwurfs
		\item Automatendarstellung, -kopplung, -vereinfachung
		\item Hardwarebeschreibungssprachen
		\item Programmierabare Schaltkreise, insbesondere FPGAs - Teil 1
		\item Programmierabare Schaltkreise, insbesondere FPGAs - Teil 2
		\item Modellierung und Simulation
		\item Zeitverhalten und Test
		\item Hochgeschwindigkeit und Verlsustleistung
		\item Anwendungsbeispiele
	\end{enumerate}

\subsection{Inhalte des Praktikums}
	\begin{enumerate}
		\item Altera Quartus-Toolchain \& Praktikumsboard DE0 mit Cyclone-3
		\item Schaltnetze und Schaltwerke
		\item Modularisierung
		\item \grqq Komplexe\grqq Anwendung: Stoppuhr
	\end{enumerate}


\section{Klassifikation von ICs}

\subsection{Zwei Sichten}
	\hfill\\Klassifizierung von integrierten Schaltkreisen (ICs) in
	\begin{itemize}
		\item Standarschaltkreis (Standard-IC) und
		\item applikationspezifische Schaltkreise (Application-Specific IC, \textbf{ASIC})
	\end{itemize}
	unter zwei Gesichtspunkten möglich:
	\begin{itemize}
		\item \textbf{Herstellungssicht}
		\item \textbf{Entwurfssicht}
	\end{itemize}

	\paragraph{Herstellungssicht:}
	\begin{itemize}
		\item Standard-IC = große Stückzahl für viele Kunden
		\item ASIC = für eine/n Kunden/Applikation speziell entwickelter und gefertigter IC mit zugeschnittener Funktionalität
	\end{itemize}

	\paragraph{Entwurfssicht:}
	\begin{itemize}
		\item Standard-IC = (Hardware-)Funktionalität kann nicht vom Anwender beeinflusst werden
		\item ASIC = vom Anwender selbst entwickelte (Hardware-)Funktionalität
	\end{itemize}

	\newpage
	\begin{center}
	\includegraphics[page=17, width=0.7\linewidth, trim=30mm 20mm 20mm 40mm, clip]{\Path/resources/Vorlesung/VLSI/01_Einfuehrung.pdf}
	\\\includegraphics[page=18, width=0.7\linewidth, trim=30mm 20mm 20mm 40mm, clip]{\Path/resources/Vorlesung/VLSI/01_Einfuehrung.pdf}
	\\\includegraphics[page=19, width=0.7\linewidth, trim=30mm 20mm 20mm 40mm, clip]{\Path/resources/Vorlesung/VLSI/01_Einfuehrung.pdf}
	\end{center}

\subsection{Anwenderprogrammierbare IC (ASIC)}
	\paragraph{Merkmale:}
	\begin{itemize}
		\item Field-Programmable $\Leftrightarrow$ feldprogrammierbar
		\item Vor Ort (im Feld) vom Anwender programmierbar
		\item Hardware ist fix. Funktionalität kann aber mittels spezieller Konfiguration \grqq programmiert\grqq werden
	\end{itemize}
	\paragraph{Anwendung:}
	\begin{itemize}
		\item Anwendungsspezifische IC bei kleinen und mittleren Stückzahlen
		\item Mehrfach neu programmierbar zwecks Optimierung und Fehlerbehebung, auch während des praktischen Einsatzes
		\item Einfache Integration eines ganzen Systems auf einem Chip
		\item Prototyping, HW-/SW-Codesign
	\end{itemize}

\subsection{Hardwareprogrammierung}
	\begin{center}
	\includegraphics[page=21, width=0.7\linewidth, trim=30mm 20mm 20mm 40mm, clip]{\Path/resources/Vorlesung/VLSI/01_Einfuehrung.pdf}
	\end{center}
%\newpage
\subsection{Programmiertechnologien}
	\begin{center}
	\includegraphics[page=22, width=0.7\linewidth, trim=30mm 20mm 20mm 40mm, clip]{\Path/resources/Vorlesung/VLSI/01_Einfuehrung.pdf}
	\\\includegraphics[page=23, width=0.7\linewidth, trim=30mm 20mm 20mm 40mm, clip]{\Path/resources/Vorlesung/VLSI/01_Einfuehrung.pdf}
	\\\includegraphics[page=24, width=0.7\linewidth, trim=30mm 20mm 20mm 40mm, clip]{\Path/resources/Vorlesung/VLSI/01_Einfuehrung.pdf}
	\\\includegraphics[page=25, width=0.7\linewidth, trim=30mm 20mm 20mm 40mm, clip]{\Path/resources/Vorlesung/VLSI/01_Einfuehrung.pdf}
	\\\includegraphics[page=26, width=0.7\linewidth, trim=30mm 20mm 20mm 40mm, clip]{\Path/resources/Vorlesung/VLSI/01_Einfuehrung.pdf}
	\\\includegraphics[page=27, width=0.7\linewidth, trim=30mm 20mm 20mm 40mm, clip]{\Path/resources/Vorlesung/VLSI/01_Einfuehrung.pdf}
	\\\includegraphics[page=28, width=0.7\linewidth, trim=30mm 20mm 20mm 40mm, clip]{\Path/resources/Vorlesung/VLSI/01_Einfuehrung.pdf}
	\end{center}

\subsection{Klassifikation Anwenderprogrammierbare IC}
	\begin{center}
	\includegraphics[page=29, width=0.7\linewidth, trim=30mm 20mm 20mm 40mm, clip]{\Path/resources/Vorlesung/VLSI/01_Einfuehrung.pdf}
	\\\includegraphics[page=30, width=0.7\linewidth, trim=30mm 20mm 20mm 40mm, clip]{\Path/resources/Vorlesung/VLSI/01_Einfuehrung.pdf}
	\\\includegraphics[page=31, width=0.7\linewidth, trim=30mm 20mm 20mm 40mm, clip]{\Path/resources/Vorlesung/VLSI/01_Einfuehrung.pdf}
	\end{center}

\newpage



\part{Entwurf eingebetteter Systeme}
\newpage
%\chapter{Einführung}
\section{Themenschwerpunkte}
	\begin{description}

		\item[Verarbeitungsleistung]\hfill
			\\\textbf{Im Vordergrund stehen:}
			\begin{itemize}
				\item Schnelle Verarbeitung auch einzelner Bits
				\item Parallelität auf: \hfill
				\begin{itemize}
					\item Bitebene
					\item Befehlseben
					\item Threadebene
					\item Prozess- und Anwendungsebene
				\end{itemize}
				\item dynamische Rekonfiguration
			\end{itemize}
			$\Rightarrow$ Erfüllung der gegebenen Anforderungen

		\item[Systemintegration]\hfill
			\\\textbf{Im Vordergrund stehen:}
			\begin{itemize}
				\item Mehrprozessorsyteme, Mehrkern-, Vielkernprozessorsysteme
				\item Mehr-Chip- / Einzel-Chip-Lösungen (System-on-a-Chip)
				\item Parallele Entwicklung von HW und SW
	(HW-/SW-Codesign)
				\item System-Prototyping (FPGA-Entwurf)
			\end{itemize}
			$\Rightarrow$ Kosteneinsparung, Entwicklungszeiteinsparung (Time-to-Market)

		\item[Verlustleistung]
			\hfill\\\textbf{Im Vordergrund stehen:}
			\begin{itemize}
				\item Verlustleistung im Standby (Akkubetrieb)
				\item Maximales Abwärmebudget
			\end{itemize}

			\hfill\\\textbf{Kenngrößen:}
			\begin{itemize}
				\item Statische und dynamische Verlustleistung
				\item MIPS pro Watt
			\end{itemize}

			$\Rightarrow$ Sowohl für eingebettete Systeme als auch für Server


		\item[Korrektheit]
			\hfill\\\textbf{Aspekte:}
			\begin{itemize}
				\item Verifikation eines Schaltkreises, simulativ / formal
				\item Profiling und Debugging unter Echtzeitbedingungen (Trace)
			\end{itemize}
			$\Rightarrow$ Fehlerfreier Erstentwurf


		\item[Fehlertoleranz]\hfill
			\\\textbf{Toleranz gegenüber:}
			\begin{itemize}
				\item Permanenten Fehlern (zeitunabhängig nach erstem Auftreten)
				\item Intermitierenden Fehlern (nur unter bestimmten Betriebsbedingungen)
				\item Transienten Fehlern (aufgrund statischer Störungen)
			\end{itemize}

			\hfill\\\textbf{Fehlererkennung und -korrektur:}
			\begin{itemize}
				\item Autonom durch Hardwarearchitektur
				\item Aus Kombination von HW und SW
			\end{itemize}

			$\Rightarrow$ Insbesondere wichtig für Sicherheitskritische, hochverfügbare und langlebige zuverlässige Systeme
			\\$\Rightarrow$ Steigende Siginifikanz mit abnehmenden Strukturgrößen (Integrationsgrad) sowie steigender Transistoranzahl (Moore's Law)
	\end{description}


\section{Inhalte der Lehrveranstaltung}
\subsection{Inhalte der Vorlesung}
	\begin{enumerate}
		\item Klassifikation von Schaltkreisen
		\item Grundlagen des Schaltkreisentwurfs
		\item Automatendarstellung, -kopplung, -vereinfachung
		\item Hardwarebeschreibungssprachen
		\item Programmierabare Schaltkreise, insbesondere FPGAs - Teil 1
		\item Programmierabare Schaltkreise, insbesondere FPGAs - Teil 2
		\item Modellierung und Simulation
		\item Zeitverhalten und Test
		\item Hochgeschwindigkeit und Verlsustleistung
		\item Anwendungsbeispiele
	\end{enumerate}

\subsection{Inhalte des Praktikums}
	\begin{enumerate}
		\item Altera Quartus-Toolchain \& Praktikumsboard DE0 mit Cyclone-3
		\item Schaltnetze und Schaltwerke
		\item Modularisierung
		\item \grqq Komplexe\grqq Anwendung: Stoppuhr
	\end{enumerate}


\section{Klassifikation von ICs}

\subsection{Zwei Sichten}
	\hfill\\Klassifizierung von integrierten Schaltkreisen (ICs) in
	\begin{itemize}
		\item Standarschaltkreis (Standard-IC) und
		\item applikationspezifische Schaltkreise (Application-Specific IC, \textbf{ASIC})
	\end{itemize}
	unter zwei Gesichtspunkten möglich:
	\begin{itemize}
		\item \textbf{Herstellungssicht}
		\item \textbf{Entwurfssicht}
	\end{itemize}

	\paragraph{Herstellungssicht:}
	\begin{itemize}
		\item Standard-IC = große Stückzahl für viele Kunden
		\item ASIC = für eine/n Kunden/Applikation speziell entwickelter und gefertigter IC mit zugeschnittener Funktionalität
	\end{itemize}

	\paragraph{Entwurfssicht:}
	\begin{itemize}
		\item Standard-IC = (Hardware-)Funktionalität kann nicht vom Anwender beeinflusst werden
		\item ASIC = vom Anwender selbst entwickelte (Hardware-)Funktionalität
	\end{itemize}

	\newpage
	\begin{center}
	\includegraphics[page=17, width=0.7\linewidth, trim=30mm 20mm 20mm 40mm, clip]{\Path/resources/Vorlesung/VLSI/01_Einfuehrung.pdf}
	\\\includegraphics[page=18, width=0.7\linewidth, trim=30mm 20mm 20mm 40mm, clip]{\Path/resources/Vorlesung/VLSI/01_Einfuehrung.pdf}
	\\\includegraphics[page=19, width=0.7\linewidth, trim=30mm 20mm 20mm 40mm, clip]{\Path/resources/Vorlesung/VLSI/01_Einfuehrung.pdf}
	\end{center}

\subsection{Anwenderprogrammierbare IC (ASIC)}
	\paragraph{Merkmale:}
	\begin{itemize}
		\item Field-Programmable $\Leftrightarrow$ feldprogrammierbar
		\item Vor Ort (im Feld) vom Anwender programmierbar
		\item Hardware ist fix. Funktionalität kann aber mittels spezieller Konfiguration \grqq programmiert\grqq werden
	\end{itemize}
	\paragraph{Anwendung:}
	\begin{itemize}
		\item Anwendungsspezifische IC bei kleinen und mittleren Stückzahlen
		\item Mehrfach neu programmierbar zwecks Optimierung und Fehlerbehebung, auch während des praktischen Einsatzes
		\item Einfache Integration eines ganzen Systems auf einem Chip
		\item Prototyping, HW-/SW-Codesign
	\end{itemize}

\subsection{Hardwareprogrammierung}
	\begin{center}
	\includegraphics[page=21, width=0.7\linewidth, trim=30mm 20mm 20mm 40mm, clip]{\Path/resources/Vorlesung/VLSI/01_Einfuehrung.pdf}
	\end{center}
%\newpage
\subsection{Programmiertechnologien}
	\begin{center}
	\includegraphics[page=22, width=0.7\linewidth, trim=30mm 20mm 20mm 40mm, clip]{\Path/resources/Vorlesung/VLSI/01_Einfuehrung.pdf}
	\\\includegraphics[page=23, width=0.7\linewidth, trim=30mm 20mm 20mm 40mm, clip]{\Path/resources/Vorlesung/VLSI/01_Einfuehrung.pdf}
	\\\includegraphics[page=24, width=0.7\linewidth, trim=30mm 20mm 20mm 40mm, clip]{\Path/resources/Vorlesung/VLSI/01_Einfuehrung.pdf}
	\\\includegraphics[page=25, width=0.7\linewidth, trim=30mm 20mm 20mm 40mm, clip]{\Path/resources/Vorlesung/VLSI/01_Einfuehrung.pdf}
	\\\includegraphics[page=26, width=0.7\linewidth, trim=30mm 20mm 20mm 40mm, clip]{\Path/resources/Vorlesung/VLSI/01_Einfuehrung.pdf}
	\\\includegraphics[page=27, width=0.7\linewidth, trim=30mm 20mm 20mm 40mm, clip]{\Path/resources/Vorlesung/VLSI/01_Einfuehrung.pdf}
	\\\includegraphics[page=28, width=0.7\linewidth, trim=30mm 20mm 20mm 40mm, clip]{\Path/resources/Vorlesung/VLSI/01_Einfuehrung.pdf}
	\end{center}

\subsection{Klassifikation Anwenderprogrammierbare IC}
	\begin{center}
	\includegraphics[page=29, width=0.7\linewidth, trim=30mm 20mm 20mm 40mm, clip]{\Path/resources/Vorlesung/VLSI/01_Einfuehrung.pdf}
	\\\includegraphics[page=30, width=0.7\linewidth, trim=30mm 20mm 20mm 40mm, clip]{\Path/resources/Vorlesung/VLSI/01_Einfuehrung.pdf}
	\\\includegraphics[page=31, width=0.7\linewidth, trim=30mm 20mm 20mm 40mm, clip]{\Path/resources/Vorlesung/VLSI/01_Einfuehrung.pdf}
	\end{center}

\newpage




\part{Parallelverarbeitung}
\newpage
%\input{../Parallelrechner/Skript/00-Parallelrechner}

\part{Appendix}
\newpage
\chapter{VLSI-Systementwurf Praktikum}
	\section{Kurze Beschreibung des Terasic DE0 Board}
	\begin{center}
	\includegraphics[page=1, width=\linewidth, trim=30mm 30mm 20mm 110mm, clip]{\Path/resources/Praktikum/VLSI/Material/Praktikumsboard/DE0_Description.pdf}
	\begin{center}
	\end{center}
	\includegraphics[page=2, width=\linewidth, trim=30mm 90mm 20mm 30mm, clip]{\Path/resources/Praktikum/VLSI/Material/Praktikumsboard/DE0_Description.pdf}
	\end{center}
	
	\def \Path {/media/Daten/Studium/Skripte/TechnischeInformatik/resources/Praktikum/VLSI/Abgabe/Protokolle}
	\begin{center}
	\includegraphics[page=1, width=\linewidth, trim=30mm 100mm 20mm 110mm, clip]{\Path/resources/Stoppuhr.pdf}
	\end{center}
	\section{Aufgabe 1 - Binär-Dekoder}
\subsection{Entwurf}
\begin{description}
	\item[Input] 4-Bit Binärzahl durch Schieberegister SW3 ... SW0
	\item[Output] 7-Segmente Darstellung einer Hexadezimalziffer (8 Einzelsignale = 7 Segmente + 1 Punkt)
\end{description}

\begin{tabular}{|c|c|c|c||c|c|c|c|c|c|c|c|c|}
\hline
\multicolumn{4}{|c|}{Input} & \multicolumn{9}{|c|}{Output} \\
\hline
\textsc{SW3} & \textsc{SW2} & \textsc{SW1} & \textsc{SW0} & 
\textsc{Hex} & \textsc{A} & \textsc{B} & \textsc{C} & \textsc{D} & \textsc{E} & \textsc{F} & \textsc{G} & \textsc{dot} \\ 
\hline
\hline
0 & 0 & 0 & 0 & 0 & 0 & 0 & 0 & 0 & 0 & 0 & 1 & 1 \\ 
0 & 0 & 0 & 1 & 1 & 1 & 0 & 0 & 1 & 1 & 1 & 1 & 1 \\
0 & 0 & 1 & 0 & 2 & 0 & 0 & 1 & 0 & 0 & 1 & 0 & 1 \\
0 & 0 & 1 & 1 & 3 & 0 & 0 & 0 & 0 & 1 & 1 & 0 & 1 \\
0 & 1 & 0 & 0 & 4 & 1 & 0 & 0 & 1 & 1 & 0 & 0 & 1 \\
0 & 1 & 0 & 1 & 5 & 0 & 1 & 0 & 0 & 1 & 0 & 0 & 1 \\
0 & 1 & 1 & 0 & 6 & 0 & 1 & 0 & 0 & 0 & 0 & 0 & 1 \\
0 & 1 & 1 & 1 & 7 & 0 & 0 & 0 & 1 & 1 & 1 & 1 & 1 \\
1 & 0 & 0 & 0 & 8 & 0 & 0 & 0 & 0 & 0 & 0 & 0 & 1 \\
1 & 0 & 0 & 1 & 9 & 0 & 0 & 0 & 0 & 1 & 0 & 0 & 1 \\
1 & 0 & 1 & 0 & A & 0 & 0 & 0 & 1 & 0 & 0 & 0 & 1 \\
1 & 0 & 1 & 1 & b & 1 & 1 & 0 & 0 & 0 & 0 & 0 & 1 \\
1 & 1 & 0 & 0 & C & 0 & 1 & 1 & 0 & 0 & 0 & 1 & 1 \\
1 & 1 & 0 & 1 & d & 1 & 0 & 0 & 0 & 0 & 1 & 0 & 1 \\
1 & 1 & 1 & 0 & E & 0 & 1 & 1 & 0 & 0 & 0 & 0 & 1 \\
1 & 1 & 1 & 1 & F & 0 & 1 & 1 & 1 & 0 & 0 & 0 & 1 \\
\hline
\end{tabular}	
\\\\
\ref{lst:01-decoder}~Decoder.vhdl Code
\subsection{Auswertung}
\paragraph{Ressourcenbedarf}
\begin{itemize} 
\item 7 Logik-Elemente
\item 12 Pins 
\end{itemize}

	\begin{center}
	\includegraphics[page=1, width=\linewidth, trim=30mm 45mm 20mm 198mm, clip]{\Path/resources/Stoppuhr.pdf}
	\end{center}
	\section{Aufgabe 2 - Zeitmessungen}
(nicht Taurus)

\begin{table}[!htb]
\caption{Original, gcc}
\begin{minipage}{.5\linewidth}
\centering
\subsection{Original, gcc, keine flags}
\begin{tabular}{|l|r|r|}
	\hline
	\textsc{Dimension} & \textsc{Runtime} & \textsc{GFLOP/s} \\
	\hline
	\hline
	32  &  0.0002s  & 0.40 \\ 
	\hline 
	64  &  0.0015s  & 0.40 \\ 
	\hline 
	96  &  0.0045s  & 0.40 \\ 
	\hline 
	128  &  0.0115s  & 0.38 \\ 
	\hline 
	160  &  0.0217s  & 0.37 \\ 
	\hline 
	192  &  0.0503s  & 0.36 \\ 
	\hline 
	224  &  0.0651s  & 0.35 \\ 
	\hline 
	256  &  0.1365s  & 0.25 \\ 
	\hline 
	320  &  0.1984s  & 0.33 \\ 
	\hline 
	384  &  0.4676s  & 0.25 \\ 
	\hline 
	448  &  0.5567s  & 0.33 \\ 
	\hline 
	512  &  1.3026s  & 0.21 \\ 
	\hline 
	640  &  2.9606s  & 0.18 \\ 
	\hline 
	768  &  5.5037s  & 0.16 \\ 
	\hline 
	896  & 8.5051s  & 0.17 \\ 
	\hline 
	1024  & 26.0022s  & 0.16 \\ 
	\hline 

\end{tabular}
\end{minipage}%
\begin{minipage}{.5\linewidth}
\centering
\subsection{Original, gcc, -O3}
\begin{tabular}{|l|r|r|}
	\hline
	\textsc{Dimension} & \textsc{Runtime} & \textsc{GFLOP/s} \\
	\hline
	\hline
	32  &  0.0000s  & 1.73 \\ 
	\hline 
	64  &  0.0003s  & 1.85 \\ 
	\hline 
	96  &  0.0009s  & 1.93 \\ 
	\hline 
	128  &  0.0031s  & 1.35 \\ 
	\hline 
	160  &  0.0055s  & 1.56 \\ 
	\hline 
	192  &  0.0104s  & 1.37 \\ 
	\hline 
	224  &  0.0159s  & 1.42 \\ 
	\hline 
	256  &  0.0339s  & 1.00 \\ 
	\hline 
	320  &  0.0537s  & 1.24 \\ 
	\hline 
	384  &  0.1077s  & 1.06 \\ 
	\hline 
	448  &  0.1524s  & 1.18 \\ 
	\hline 
	512  &  0.3500s  & 0.78 \\ 
	\hline 
	640  &  1.6626s  & 0.33 \\ 
	\hline 
	768  &  3.2088s  & 0.29 \\ 
	\hline 
	896  & 5.1806s  & 0.28 \\ 
	\hline 
	1024  & 7.5068s  & 0.29 \\ 
	\hline 

\end{tabular}

\end{minipage}%
\end{table}


\begin{table}[!htb]
\caption{Mit Optimierungen, gcc}
\begin{minipage}{.5\linewidth}
\centering
\subsection{Mit Optimierungen, gcc, keine flags}
\begin{tabular}{|l|r|r|}
	\hline
	\textsc{Dimension} & \textsc{Runtime} & \textsc{GFLOP/s} \\
	\hline
	\hline
	32  &  0.0001s  & 0.61 \\ 
	\hline 
	64  &  0.0009s  & 0.61 \\ 
	\hline 
	96  &  0.0026s  & 0.70 \\ 
	\hline 
	128  &  0.0058s  & 0.71 \\ 
	\hline 
	160  &  0.0115s  & 0.71 \\ 
	\hline 
	192  &  0.0198s  & 0.71 \\ 
	\hline 
	224  &  0.0313s  & 0.72 \\ 
	\hline 
	256  &  0.0465s  & 0.72 \\ 
	\hline 
	320  &  0.0902s  & 0.73 \\ 
	\hline 
	384  &  0.1552s  & 0.73 \\ 
	\hline 
	448  &  0.2451s  & 0.73 \\ 
	\hline 
	512  &  0.3632s  & 0.74 \\ 
	\hline 
	640  &  0.7120s  & 0.74 \\ 
	\hline 
	768  &  1.2261s  & 0.74 \\ 
	\hline 
	896  &  1.9538s  & 0.74 \\ 
	\hline 
	1024  &  2.9417s  & 0.73 \\ 
	\hline 
\end{tabular}
\end{minipage}%
\begin{minipage}{.5\linewidth}
\centering
\subsection{Mit Optimierungen, gcc, -O3}
\begin{tabular}{|l|r|r|}
	\hline
	\textsc{Dimension} & \textsc{Runtime} & \textsc{GFLOP/s} \\
	\hline
	\hline
	32  &  0.0000s  & 2.12 \\ 
	\hline 
	64  &  0.0002s  & 2.24 \\ 
	\hline 
	96  &  0.0006s  & 2.50 \\ 
	\hline 
	128  &  0.0015s  & 2.58 \\ 
	\hline 
	160  &  0.0031s  & 2.56 \\ 
	\hline 
	192  &  0.0052s  & 2.65 \\ 
	\hline 
	224  &  0.0081s  & 2.67 \\ 
	\hline 
	256  &  0.0120s  & 2.75 \\ 
	\hline 
	320  &  0.0239s  & 2.72 \\ 
	\hline 
	384  &  0.0408s  & 2.76 \\ 
	\hline 
	448  &  0.0647s  & 2.77 \\ 
	\hline 
	512  &  0.0941s  & 2.86 \\ 
	\hline 
	640  &  0.1886s  & 2.78 \\ 
	\hline 
	768  &  0.3226s  & 2.81 \\ 
	\hline 
	896  &  0.5146s  & 2.80 \\ 
	\hline 
	1024  &  0.7530s  & 2.85 \\ 
	\hline 
\end{tabular}
\end{minipage}%
\end{table}


\begin{table}[!htb]
\caption{Original, icc}
\begin{minipage}{.5\linewidth}
\centering
\subsection{Original, icc, keine flags}
\begin{tabular}{|l|r|r|}
	\hline
	\textsc{Dimension} & \textsc{Runtime} & \textsc{GFLOP/s} \\
	\hline
	\hline
	32  &  0.0000s  & 1.52 \\ 
	\hline 
	64  &  0.0001s  & 3.67 \\ 
	\hline 
	96  &  0.0003s  & 4.30 \\ 
	\hline 
	128  &  0.0008s  & 4.88 \\ 
	\hline 
	160  &  0.0018s  & 4.87 \\ 
	\hline 
	192  &  0.0030s  & 4.90 \\ 
	\hline 
	224  &  0.0046s  & 4.73 \\ 
	\hline 
	256  &  0.0072s  & 4.73 \\ 
	\hline 
	320  &  0.0125s  & 4.99 \\ 
	\hline 
	384  &  0.0228s  & 5.02 \\ 
	\hline 
	448  &  0.0346s  & 5.17 \\ 
	\hline 
	512  &  0.0533s  & 5.02 \\ 
	\hline 
	640  &  0.1020s  & 5.13 \\ 
	\hline 
	768  &  0.1739s  & 5.15 \\ 
	\hline 
	896  &  0.2743s  & 5.25 \\ 
	\hline 
	1024  &  0.4137s  & 5.19 \\ 
	\hline 
\end{tabular}
\end{minipage}%
\begin{minipage}{.5\linewidth}
\centering
\subsection{Original, icc, -O3}
\begin{tabular}{|l|r|r|}
	\hline
	\textsc{Dimension} & \textsc{Runtime} & \textsc{GFLOP/s} \\
	\hline
	\hline
	32  &  0.0000s  & 1.98 \\ 
	\hline 
	64  &  0.0001s  & 3.69 \\ 
	\hline 
	96  &  0.0004s  & 4.00 \\ 
	\hline 
	128  &  0.0009s  & 4.51 \\ 
	\hline 
	160  &  0.0016s  & 4.67 \\ 
	\hline 
	192  &  0.0031s  & 4.60 \\ 
	\hline 
	224  &  0.0051s  & 4.43 \\ 
	\hline 
	256  &  0.0076s  & 4.40 \\ 
	\hline 
	320  &  0.0144s  & 4.57 \\ 
	\hline 
	384  &  0.0245s  & 4.66 \\ 
	\hline 
	448  &  0.0386s  & 4.67 \\ 
	\hline 
	512  &  0.0582s  & 4.60 \\ 
	\hline 
	640  &  0.1117s  & 4.70 \\ 
	\hline 
	768  &  0.1937s  & 4.69 \\ 
	\hline 
	896  &  0.3035s  & 4.75 \\ 
	\hline 
	1024  &  0.4552s  & 4.73 \\ 
	\hline 
\end{tabular}
\end{minipage}%
\end{table}




\begin{table}[!htb]
\caption{Mit Optimierungen, icc}
\begin{minipage}{.5\linewidth}
\centering
\subsection{Mit Optimierungen, icc, keine flags}
\begin{tabular}{|l|r|r|}
	\hline
	\textsc{Dimension} & \textsc{Runtime} & \textsc{GFLOP/s} \\
	\hline
	\hline
	32  &  0.0000s  & 2.99 \\ 
	\hline 
	64  &  0.0001s  & 5.09 \\ 
	\hline 
	96  &  0.0003s  & 5.41 \\ 
	\hline 
	128  &  0.0007s  & 5.98 \\ 
	\hline 
	160  &  0.0014s  & 5.95 \\ 
	\hline 
	192  &  0.0024s  & 5.65 \\ 
	\hline 
	224  &  0.0041s  & 5.41 \\ 
	\hline 
	256  &  0.0062s  & 5.29 \\ 
	\hline 
	320  &  0.0121s  & 5.45 \\ 
	\hline 
	384  &  0.0202s  & 5.52 \\ 
	\hline 
	448  &  0.0327s  & 5.57 \\ 
	\hline 
	512  &  0.0481s  & 5.61 \\ 
	\hline 
	640  &  0.0928s  & 5.63 \\ 
	\hline 
	768  &  0.1606s  & 5.65 \\ 
	\hline 
	896  &  0.2522s  & 5.70 \\ 
	\hline 
	1024  &  0.3757s  & 5.72 \\ 
	\hline 
\end{tabular}
\end{minipage}%
\begin{minipage}{.5\linewidth}
\centering
\subsection{Mit Optimierungen, icc, -O3}
\begin{tabular}{|l|r|r|}
	\hline
	\textsc{Dimension} & \textsc{Runtime} & \textsc{GFLOP/s} \\
	\hline
	\hline
	32  &  0.0000s  & 2.72 \\ 
	\hline 
	64  &  0.0001s  & 5.07 \\ 
	\hline 
	96  &  0.0003s  & 5.34 \\ 
	\hline 
	128  &  0.0006s  & 5.67 \\ 
	\hline 
	160  &  0.0016s  & 5.60 \\ 
	\hline 
	192  &  0.0024s  & 5.78 \\ 
	\hline 
	224  &  0.0046s  & 5.35 \\ 
	\hline 
	256  &  0.0062s  & 5.32 \\ 
	\hline 
	320  &  0.0122s  & 5.44 \\ 
	\hline 
	384  &  0.0202s  & 5.51 \\ 
	\hline 
	448  &  0.0322s  & 5.58 \\ 
	\hline 
	512  &  0.0482s  & 5.57 \\ 
	\hline 
	640  &  0.0950s  & 5.52 \\ 
	\hline 
	768  &  0.1650s  & 5.50 \\ 
	\hline 
	896  &  0.2561s  & 5.54 \\ 
	\hline 
	1024  &  0.3841s  & 5.58 \\ 
	\hline 
\end{tabular}
\end{minipage}%
\end{table}

	\begin{center}
	\includegraphics[page=2, width=\linewidth, trim=30mm 130mm 20mm 30mm, clip]{\Path/resources/Stoppuhr.pdf}
	\end{center}
	\section{Aufgabe 3 - Compiler-Flags}

\subsection{-O3}
Bei Verwendung des gcc-Compilers bringt dieser Flag eine Verbesserung der Ausführungszeit vom Faktor 4 mit sich. Allerdings verschlechtert er die Ausführungszeit beim icc.
\subsection{-floop-interchange}
Führt eine Vertauschung von Schleifen aus, ähnlich wie die Code-Optimierung.
\subsection{-funroll-loops}
Nimmt Schleifen auseinander, deren Schritte durch den Compiler vor der Ausführung bestimmt werden können.



	\begin{center}
	\includegraphics[page=3, width=\linewidth, trim=30mm 100mm 20mm 30mm, clip]{\Path/resources/Stoppuhr.pdf}
	\end{center}
	\section{Aufgabe 4 - Entprell-Automat}
\subsection{Entwurf}
	\paragraph{State-Machine-Charts}\hfill \\

	\paragraph{LED} enthält die Komponente Entprellung und verbindet die Ein- und Ausgangssignale. (\ref{lst:04-led}~LED.vhdl-Code) \\
	\includegraphics[width=0.7\textwidth]{\Path/resources/04-led.png}
	
	\paragraph{Entprellung} enthält die Komponente Zaehler, der bei der Veränderung des Eingangsignals gestartet wird und für 3ms weitere Änderungen ignoriert. (\ref{lst:04-entprellung}~Entprellung.vhdl-Code) \\
	\includegraphics[width=0.8\textwidth]{\Path/resources/04-entpreller.png}


	\paragraph{Zaehler} implementiert einen Zähler, der durch ein Signal definierte Schritte zählt. Ausgegeben wird der aktuelle Zustand des Zählers. Eingegeben ein Reset-Signal. (\ref{lst:04-zaehler}~Zaehler.vhdl-Code) 
	\paragraph{Kopplung} \hfill\\
	Es wird eine synchrone Automatenkopplung über die Ausgangssignale mit einem Taktsignal verwendet.


\subsection{Auswertung}
	\paragraph{Ressourcenbedarf}
	\begin{itemize} 
	\item 86 Logik-Elemente
	\item davon 79 dedizierte Logik-Elemente
	\item 44 Register
	\item 3 Pins 
	\item maximale Taktfrequenz von 178 MHz
	\end{itemize}

	\begin{center}
	\includegraphics[page=4, width=\linewidth, trim=30mm 130mm 20mm 30mm, clip]{\Path/resources/Stoppuhr.pdf}
	\end{center}
	\section{Aufgabe 5 - HALLO-Anzeige}
\subsection{Entwurf}
	\begin{description}
	\item[zu a)] Es müssen 5 Zeichen kodiert werden (H, A, L, O, Leerzeichen).\\
			\\\(ld\, 5 = 3\)\\\\
			Daher werden für eine Binärkodierung mindestens 3 Bits benötigt.\\\\
			\begin{tabular}{|c|c|c||c|c|c|c|c|c|c|c|c|}
		\hline
		\multicolumn{3}{|c||}{Input} & \multicolumn{9}{|c|}{Output} \\
		\hline
		\textsc{BIT2} & \textsc{BIT1} & \textsc{BIT0} & 
		\textsc{CHAR} & \textsc{A} & \textsc{B} & \textsc{C} & \textsc{D} & \textsc{E} & \textsc{F} & \textsc{G} & \textsc{dot} \\ 
		\hline
		\hline
		0 & 0 & 0 &   & 1 & 1 & 1 & 1 & 1 & 1 & 1 & 1 \\ 
		0 & 0 & 1 & H & 1 & 0 & 0 & 1 & 0 & 0 & 0 & 1 \\
		0 & 1 & 0 & A & 0 & 0 & 0 & 1 & 0 & 0 & 0 & 1 \\
		0 & 1 & 1 & L & 1 & 1 & 1 & 0 & 0 & 0 & 1 & 1 \\
		1 & 0 & 0 & O & 0 & 0 & 0 & 0 & 0 & 0 & 1 & 1 \\
		\hline
	\end{tabular}
	\item[b)] Für das Schieberegister ist der Zählerzustand ein Enable-Signal
	\item[c)]  (\ref{lst:05-hallo}~Hallo.vhdl-Code) 
	\end{description}

	

\subsection{Auswertung}
	\paragraph{Ressourcenbedarf}
	\begin{itemize} 
	\item 73 Logik-Elemente
	\item 61 Register
	\item 34 Pins 
	\item maximale Taktfrequenz von 262 MHz
	\end{itemize}

	\begin{center}
	\includegraphics[page=4, width=\linewidth, trim=30mm 25mm 20mm 175mm, clip]{\Path/resources/Stoppuhr.pdf}
	\includegraphics[page=5, width=\linewidth, trim=30mm 200mm 20mm 30mm, clip]{\Path/resources/Stoppuhr.pdf}
	\end{center}
	\section{Aufgabe 6 - Stoppuhr}
\subsection{Entwurf}
\includegraphics[width=1\textwidth]{\Path/resources/06-stopuhr.png}
	\newpage
	\paragraph{State-Machine-Charts}\hfill \\

		\includegraphics[width=0.8\textwidth]{\Path/resources/06-Entprellautomat.png}
		

		\includegraphics[width=0.7\textwidth]{\Path/resources/06-Zaehler.png}

		\includegraphics[width=0.7\textwidth]{\Path/resources/06-Stoppuhr.png}
 
	\paragraph{Kopplung} \hfill\\
	Es wird eine synchrone Automatenkopplung über die Ausgangssignale mit einem Taktsignal verwendet.
\subsection{Auswertung}
	\paragraph{Ressourcenbedarf}
	\begin{itemize} 
	\item 163 Logik-Elemente
	\item davon 121 dedizierte Logik-Elemente
	\item 35 Pins 
	\item maximale Taktfrequenz von 225 MHz
	\end{itemize}

	\section{Anhang}
\subsection{Circuit}
\label{lst:circuit}
\lstinputlisting[caption={Circuit.h}, language=C++]{\Path/../Simulator/src/main/Circuit.h}

\subsection{Simulator}
\label{lst:simulator}
\lstinputlisting[caption={Simulator.h}, language=C++]{\Path/../Simulator/src/main/Simulator.h}

\subsection{Solver}
\label{lst:solver}
\lstinputlisting[caption={Solver.h}, language=C++]{\Path/../Simulator/src/main/Solver.h}

\subsection{Parsers}
\label{lst:parsers}
\lstinputlisting[caption={Parsers.h}, language=C++]{\Path/../Simulator/src/parser/Parsers.h}

\subsection{Parser}
\label{lst:parser}
\lstinputlisting[caption={Parser.h}, language=C++]{\Path/../Simulator/src/parser/Parser.h}

\subsection{BENCH}
\label{lst:bench}
\lstinputlisting[caption={BENCH.h}, language=C++]{\Path/../Simulator/src/parser/BENCH.h}

\subsection{Gates}
\label{lst:gates}
\lstinputlisting[caption={Gates.h}, language=C++]{\Path/../Simulator/src/gates/Gates.h}

\subsection{Gate}
\label{lst:gate}
\lstinputlisting[caption={Gate.h}, language=C++]{\Path/../Simulator/src/gates/Gate.h}

\subsection{Input}
\label{lst:input}
\lstinputlisting[caption={Input.h}, language=C++]{\Path/../Simulator/src/gates/Input.h}

\subsection{DFF}
\label{lst:dff}
\lstinputlisting[caption={DFF.h}, language=C++]{\Path/../Simulator/src/gates/DFF.h}

\subsection{Output}
\label{lst:output}
\lstinputlisting[caption={Output.h}, language=C++]{\Path/../Simulator/src/gates/Output.h}
\newpage


	\def \Path {/media/Daten/Studium/Skripte/TechnischeInformatik}


\chapter{Entwurf eingebetteter Systeme: Schaltkreisvalidation}
	
	\def \Path {/media/Daten/Studium/Skripte/TechnischeInformatik/resources/Praktikum/Mikrorechner/Protokolle}
	\section{Programm}

\subsection{Entwurf}
\subsubsection{Circuit}
Das Grundgerüst des Programms bildet die \code{Circuit}-Klasse, die einen Schaltkreis und seine Funktionalitäten implementiert. Die Eingangsbelegung des Schaltkreises wird durch eine Zustandsvariable \code{unsigned int state} repräsentiert, die einen Wert zwischen 0 und $2^N$ annehmen kann, wobei $N$ die Anzahl der Eingänge ist. 
\\\\
Die Eingänge (\code{Input}-Objekte), Gatter (\code{Gate}-Objekte) und Ausgänge (\code{Output}-Objekte) werden in jeweils eigenen Maps abgelegt, deren Keys vom Typ \code{std::string} die Bezeichner der jeweiligen Objekte sind. Alle Gatter sind als Klassen implementiert, die von der abstrakten Klasse \code{Gate} erben. Bei manchen Gattern - wie AND, Or, usw. - können beliebig viele Eingänge definiert werden. Ein Gattereingang ist ein Zeiger, auf entweder den Wert eines \code{Input}-Objektes oder auf den Ausgang eines anderen \code{Gate}-Objektes. Die Ein- und Ausgänge der Gatter sind vom Typ \code{bool}.
\\\\
Bei der Simulation eines Schaltkreises wird über den \code{state} der \code{Circuit}-Klasse iteriert und dessen Wert in Binärdarstellung auf die Eingänge abgebildet. Nach jeder neuen Eingangsbelegung müssen die Gatter durchlaufen werden, bis das neue Signal die Ausgänge erreicht.

\subsubsection{Parser}
Zum Parsen wird die Klasse \code{Parser} als Grundgerüst zur Verfügung gestellt, wobei die genauere Implementierung für unterschiedliche Benchmarks in von dieser Klasse erbenden Klassen verschoben wurde. Für das Benchmark BENCH ist eine solche gleichnamige Klasse vorhanden. 
\\\\
Ein Parser kann Schaltkreise aus Dateien lesen und diese anschliesend als \code{Circuit}-Objekte zurückgeben.
\\\\
Der Parser für das BENCH-Format geht davon aus, dass in der einzulesenden Datei zuerst alle Eingänge (INPUT), dann die Ausgänge (OUTPUT), danach die ggf. vorhandenen FlipFlops (DFF) und zuletzt die restlichen Gatter definiert werden. 
\\\\
Um einen Schaltkreis aus einer Datei zu Parsen, müssen dem Programm die Argumente \\\code{-p BENCHMARK DATEI} mitgeteilt werden.

\subsubsection{Simulator}
Die Klasse \code{Simulator} enthält eine Liste von geparsten \code{Circuit}-Objekten, die simuliert werden können. Dazu wird wieder über den Zustand der Schaltkreise iteriert.
\\\\
Um zum Beispiel zwei Schaltkreise aus Dateien zu parsen und zu Simulieren müssen folgende Argumente übergeben werden:\\
\code{-p BENCHMARK1 DATEI1 -p BENCHMARK2 DATEI2 -s}

	
	\newpage
	\subsection{Äquivalenzprüfung duch Simulation}

\subsubsection{Voraussetzungen}
\begin{itemize}
	\item Schaltkreise müssen gleichviele Eingänge bzw. Ausgänge besitzen
	\item falls in einem Schaltkreis ein Eingang bzw. Ausgang vorkommt, muss ein gleichnamiger Eingang bzw. Ausgang auch in den anderen Schaltkreisen vorkommen
	\item es können beliebig viele Schaltkreise verglichen werden, falls einer nicht äquivalent mit einem anderen ist, wird \code{false} zurückgegeben
	\item ist kein Schaltkreis definiert, wird \code{false} zurückgegeben
	\item ist nur ein Schaltkreis definiert, wird \code{true} zurückgegeben

\end{itemize}

\subsubsection{Vorgehen}
\begin{enumerate}
	\item über mögliche Zustände ($2^N$ mit N ... Anzahl der Eingänge) iterieren

	\begin{enumerate}
		\item durch Schaltkreise iterieren
		\begin{enumerate}
			\item Zustand auf Eingänge abbilden
			\item Schaltkreis traversieren und Ausgangsbelegung ermitteln
		\end{enumerate}
		\item Belegungen mit denen des vorangegangenen Schaltkreises Vergleichen
		\item Bei unterschiedlicher Belegung wird \code{false} zurückgegeben
		\item Sonst wird die Iteration über die Zustände fortgesetzt
	\end{enumerate}
	
	\item wurde über alle $2^N$ Zustände iteriert und keine Varianz der Ausgangsbelegung bei den Schaltkreisen festgestellt, wird \code{true} zurückgegeben
\end{enumerate}


\subsubsection{Implementierung}
Für alle Circuits werden die möglichen Inputs berechnet und die erhaltenen Outputs miteinander verglichen. Sollte dabei ein Circuit enthalten sein, der abweicht, wird false zurückgegeben (Circuits sind nicht äquivalent). 

	
	\newpage
	\subsection{Äquivalenzprüfung durch SAT-Solver}

\subsubsection{Voraussetzungen}
\begin{itemize}
	\item Schaltkreise müssen gleichviele Eingänge bzw. Ausgänge besitzen
	\item falls in einem Schaltkreis ein Eingang bzw. Ausgang vorkommt, muss ein gleichnamiger Eingang bzw. Ausgang auch in den anderen Schaltkreisen vorkommen
	\item es können beliebig viele Schaltkreise verglichen werden, falls einer nicht äquivalent mit einem anderen ist, wird \code{false} zurückgegeben
	\item ist kein Schaltkreis definiert, wird \code{false} zurückgegeben
	\item ist nur ein Schaltkreis definiert, wird \code{true} zurückgegeben

\end{itemize}


\subsubsection{Vorgehen}
\begin{enumerate}
	\item es wird über alle Schaltkreise iteriert
	
	\begin{enumerate}
		\item die KNFs aller Gatter werden ermittelt
		\item falls der Schaltkreis nicht der erste ist, werden seine Eingänge mit denen des ersten Schaltkreises verbunden d.h. die Identität wird duch eine KNF dargestellt
		\item falls der Schaltkreis nicht der erste ist, werden seine Ausgänge XOR mit denen des ersten Schaltkreises verbunden d.h. als KNF dargestellt
	\end{enumerate}
	
	\item alle Ausgänge der XOR-Verbindungen werden OR verküpft, d.h. es wird eine Klausel gebildet, die alle diese Ausgänge enthält
	\subitem lässt sich eine dieser Variablen auf \code{true} abbilden (die These ist SATISFIABLE), dann sind die definierten Schaltkreise nicht äquivalent 
	\item gesamte KNF wird in Datei geschrieben und der SATsolver gestartet
\end{enumerate}


\subsubsection{Implementierung}
Der verwendete SAT-Solver ist miniSAT in der Version 1.14 als Binary. Die generierten KNF-Klauseln werden in eine Datei mit dem Namen cnf geschrieben und anschliesend vom SAT-Solver gelesen. Die jeweiligen Gatter-Klassen sind so implementiert, dass sie ihre Funktion als KNF darstellen können. Die Klasse \code{Solver} enthält einen Zähler für die Benennung der KNF-Variablen.

	
	\newpage
	\chapter{Versuche}

\section{1. Versuch}

\begin{tabular}{|l|c|c|}
	\hline
	\textsc{Benchmarks} & tests/01-test/s344.bench & tests/01-test/s349.bench \\
	\hline
	\hline
	\textsc{Parser} & BENCH & BENCH \\
	\hline
	\textsc{Inputs} & 9 & 9 \\
	\hline
	\textsc{DFFs} & 15 & 15 \\
	\hline
	\textsc{Outputs} & 11 & 11 \\
	\hline
	\textsc{Gatter} & 160 & 161 \\		
	\hline
	\textsc{Zeit in Sekunden} & 0.105102 &  0.103989 \\ 
	\hline
	\hline
	\textsc{Äquivalenzcheck durch SAT-Solving} & \multicolumn{2}{|c|}{äquivalent} \\
	\hline
	\textsc{Zeit in Sekunden} & \multicolumn{2}{|c|}{0.196012} \\
	\hline
	\hline
	\textsc{Äquivalenzcheck durch Simmulierung} & \multicolumn{2}{|c|}{nicht ausgeführt} \\
	\hline
	\textsc{Zeit in Sekunden} & \multicolumn{2}{|c|}{ewig} \\
	\hline
\end{tabular}
\\\\ 
Zu lange Ausführungszeiten für den Simulator.

\section{2. Versuch}

\begin{tabular}{|l|c|c|}
	\hline
	\textsc{Benchmarks} & tests/02-test/s298.bench & tests/02-test/s298.bench \\
	\hline
	\hline
	\textsc{Parser} & BENCH & BENCH \\
	\hline
	\textsc{Inputs} & 3 & 3 \\
	\hline
	\textsc{DFFs} & 14 & 14 \\
	\hline
	\textsc{Outputs} & 6 & 6 \\
	\hline
	\textsc{Gatter} & 119 & 119 \\		
	\hline
	\textsc{Zeit in Sekunden} & 0.079251 &  0.05012 \\ 
	\hline
	\hline
	\textsc{Äquivalenzcheck durch SAT-Solving} & \multicolumn{2}{|c|}{äquivalent} \\
	\hline
	\textsc{Zeit in Sekunden} & \multicolumn{2}{|c|}{0.096006} \\
	\hline
	\hline
	\textsc{Äquivalenzcheck durch Simmulierung} & \multicolumn{2}{|c|}{äquivalent} \\
	\hline
	\textsc{Zeit in Sekunden} & \multicolumn{2}{|c|}{2034.31} \\
	\hline
\end{tabular}

\section{3. Versuch}

\begin{tabular}{|l|c|c|}
	\hline
	\textsc{Benchmarks} & tests/03-test/s298.bench & tests/03-test/s298a.bench \\
	\hline
	\hline
	\textsc{Parser} & BENCH & BENCH \\
	\hline
	\textsc{Inputs} & 3 & 3 \\
	\hline
	\textsc{DFFs} & 14 & 14 \\
	\hline
	\textsc{Outputs} & 6 & 6 \\
	\hline
	\textsc{Gatter} & 119 & 119 \\		
	\hline
	\textsc{Zeit in Sekunden} & 0.048399 &  0.047988 \\ 
	\hline
	\hline
	\textsc{Äquivalenzcheck durch SAT-Solving} & \multicolumn{2}{|c|}{nicht äquivalent} \\
	\hline
	\textsc{Zeit in Sekunden} & \multicolumn{2}{|c|}{0.100006} \\
	\hline
	\hline
	\textsc{Äquivalenzcheck durch Simmulierung} & \multicolumn{2}{|c|}{nicht äquivalent} \\
	\hline
	\textsc{Zeit in Sekunden} & \multicolumn{2}{|c|}{0.016245} \\
	\hline
\end{tabular}

\section{4. Versuch}

\begin{tabular}{|l|c|c|}
	\hline
	\textsc{Benchmarks} & tests/04-test/c17.bench & tests/04-test/c17a.bench \\
	\hline
	\hline
	\textsc{Parser} & BENCH & BENCH \\
	\hline
	\textsc{Inputs} & 5 & 5 \\
	\hline
	\textsc{DFFs} & 0 & 0 \\
	\hline
	\textsc{Outputs} & 2 & 2 \\
	\hline
	\textsc{Gatter} & 6 & 6 \\		
	\hline
	\textsc{Zeit in Sekunden} & 0.002447 & 0.00174 \\ 
	\hline
	\hline
	\textsc{Äquivalenzcheck durch SAT-Solving} & \multicolumn{2}{|c|}{nicht äquivalent} \\
	\hline
	\textsc{Zeit in Sekunden} & \multicolumn{2}{|c|}{0.004} \\
	\hline
	\hline
	\textsc{Äquivalenzcheck durch Simmulierung} & \multicolumn{2}{|c|}{nicht äquivalent} \\
	\hline
	\textsc{Zeit in Sekunden} & \multicolumn{2}{|c|}{0.000122} \\
	\hline
\end{tabular}

\section{5. Versuch}

\begin{tabular}{|l|c|c|}
	\hline
	\textsc{Benchmarks} & tests/05-test/c17.bench & tests/05-test/c17a.bench \\
	\hline
	\hline
	\textsc{Parser} & BENCH & BENCH \\
	\hline
	\textsc{Inputs} & 5 & 5 \\
	\hline
	\textsc{DFFs} & 0 & 0 \\
	\hline
	\textsc{Outputs} & 2 & 2 \\
	\hline
	\textsc{Gatter} & 6 & 7 \\		
	\hline
	\textsc{Zeit in Sekunden} & 0.002533 & 0.001905 \\ 
	\hline
	\hline
	\textsc{Äquivalenzcheck durch SAT-Solving} & \multicolumn{2}{|c|}{äquivalent} \\
	\hline
	\textsc{Zeit in Sekunden} & \multicolumn{2}{|c|}{0} \\
	\hline
	\hline
	\textsc{Äquivalenzcheck durch Simmulierung} & \multicolumn{2}{|c|}{äquivalent} \\
	\hline
	\textsc{Zeit in Sekunden} & \multicolumn{2}{|c|}{ 0.001855} \\
	\hline
\end{tabular} \\\\
Veränderung: ein NAND-Gatter durch ein AND und ein NOT ersetzt. Ergebnis korrekt.


\section{6. Versuch}

\begin{tabular}{|l|c|c|}
	\hline
	\textsc{Benchmarks} & tests/06-test/c6288.bench & tests/05-test/c6288a.bench \\
	\hline
	\hline
	\textsc{Parser} & BENCH & BENCH \\
	\hline
	\textsc{Inputs} & 32 & 32 \\
	\hline
	\textsc{DFFs} & 0 & 0 \\
	\hline
	\textsc{Outputs} & 32 & 32 \\
	\hline
	\textsc{Gatter} & 2416 & 2416 \\		
	\hline
	\textsc{Zeit in Sekunden} & 8.64887 & 8.4408 \\ 
	\hline
	\hline
	\textsc{Äquivalenzcheck durch SAT-Solving} & \multicolumn{2}{|c|}{nicht äquivalent} \\
	\hline
	\textsc{Zeit in Sekunden} & \multicolumn{2}{|c|}{17.3411} \\
	\hline
	\hline
	\textsc{Äquivalenzcheck durch Simmulierung} & \multicolumn{2}{|c|}{nicht äquivalent} \\
	\hline
	\textsc{Zeit in Sekunden} & \multicolumn{2}{|c|}{6.16808} \\
	\hline
\end{tabular} \\\\
Veränderung: 2 Gatter geändert. 

	
	\newpage
	\section{Anhang}
\subsection{Circuit}
\label{lst:circuit}
\lstinputlisting[caption={Circuit.h}, language=C++]{\Path/../Simulator/src/main/Circuit.h}

\subsection{Simulator}
\label{lst:simulator}
\lstinputlisting[caption={Simulator.h}, language=C++]{\Path/../Simulator/src/main/Simulator.h}

\subsection{Solver}
\label{lst:solver}
\lstinputlisting[caption={Solver.h}, language=C++]{\Path/../Simulator/src/main/Solver.h}

\subsection{Parsers}
\label{lst:parsers}
\lstinputlisting[caption={Parsers.h}, language=C++]{\Path/../Simulator/src/parser/Parsers.h}

\subsection{Parser}
\label{lst:parser}
\lstinputlisting[caption={Parser.h}, language=C++]{\Path/../Simulator/src/parser/Parser.h}

\subsection{BENCH}
\label{lst:bench}
\lstinputlisting[caption={BENCH.h}, language=C++]{\Path/../Simulator/src/parser/BENCH.h}

\subsection{Gates}
\label{lst:gates}
\lstinputlisting[caption={Gates.h}, language=C++]{\Path/../Simulator/src/gates/Gates.h}

\subsection{Gate}
\label{lst:gate}
\lstinputlisting[caption={Gate.h}, language=C++]{\Path/../Simulator/src/gates/Gate.h}

\subsection{Input}
\label{lst:input}
\lstinputlisting[caption={Input.h}, language=C++]{\Path/../Simulator/src/gates/Input.h}

\subsection{DFF}
\label{lst:dff}
\lstinputlisting[caption={DFF.h}, language=C++]{\Path/../Simulator/src/gates/DFF.h}

\subsection{Output}
\label{lst:output}
\lstinputlisting[caption={Output.h}, language=C++]{\Path/../Simulator/src/gates/Output.h}
\newpage


	\def \Path {/media/Daten/Studium/Skripte/TechnischeInformatik}

\end{document}
